\documentclass[10pt,a4paper]{article}
\usepackage[margin=2.2cm]{geometry}
\usepackage{fontspec}
\setmonofont{DejaVu Sans Mono}[Scale=0.84]
\newfontfamily\unifont{DejaVu Serif}[Scale=0.9]
\usepackage{amsmath,amssymb,mathtools}
\setcounter{MaxMatrixCols}{20}
\usepackage{longtable,booktabs,array}
\usepackage{fancyvrb}
\usepackage[dvipsnames]{xcolor}
\usepackage[colorlinks=true,linkcolor=NavyBlue,urlcolor=NavyBlue]{hyperref}
\usepackage{microtype}
\usepackage{needspace}
\usepackage{tocloft}\setlength{\cftbeforesecskip}{0.45em}
\setlength{\parindent}{0pt}\setlength{\parskip}{0.55em}
\setlength{\emergencystretch}{8em}
\makeatletter
\renewcommand\section{\@startsection{section}{1}{\z@}{-3.5ex \@plus -1ex \@minus -.2ex}{2.3ex \@plus.2ex}{\normalfont\Large\bfseries\raggedright}}
\renewcommand\subsection{\@startsection{subsection}{2}{\z@}{-3.25ex\@plus -1ex \@minus -.2ex}{1.5ex \@plus .2ex}{\normalfont\large\bfseries\raggedright}}
\renewcommand\subsubsection{\@startsection{subsubsection}{3}{\z@}{-3.25ex\@plus -1ex \@minus -.2ex}{1.5ex \@plus .2ex}{\normalfont\normalsize\bfseries\raggedright}}
\makeatother
\renewcommand{\arraystretch}{1.15}
\allowdisplaybreaks
\begin{document}
\title{The auxiliary torus of 7+1 spacetime\\[0.4em]\large What it is, how it is described, and what it does — a guide for students, with worked examples checked in sections 25 and 26 of the notebook 7+1.nb and by independent Python programs}
\author{Repository: github.com/once-ere/7plus1 — prepared for Nash\\ Notebook: 7+1.nb (Wolfram Language 15.0.1), evaluated Wed 28 Oct 2026 08:57:47\\ Document generated by tools/make\_docs.py on 2026-10-28}
\date{}
\maketitle
\begin{abstract}
In the 7+1 model the four auxiliary space-like coordinates $x_4,\dots,x_7$ are periodic: they form a four-dimensional torus, $T^4=\mathbb R^4/(L_c\mathbb Z)^4$, the AUXILIARY TORUS. This document defines it precisely — as a quotient space, with its metric, its size $b$ ($b^8=\det(-g_{kl})$), its proper volume $V_\perp=b^4L_c^4$ and its symmetries — and explains what it does: fields on it are Fourier series, and their momenta along the torus appear in ordinary spacetime as the masses of a Kaluza–Klein tower; the number of levels below a mass $\mu$ is $V_\perp\mu^4/(32\pi^2)$; the volume fixes the four-dimensional Newton constant, $8\pi G_4=\kappa/V_\perp$; today ($M_8=10$ TeV) the torus has the proper size $\ell_\perp=3.0792\times10^{-13}$ m. A central worked example shows how gravity changes from the eight-dimensional $1/r^5$ law to Newton’s $1/r$ law: the sum over the images of a mass on the torus equals the sum over its Kaluza–Klein modes (Poisson summation), and the long-range force is $\frac53G_4$ — the graviton plus the radion — until the stabilizing potential gives the radion its mass. The document follows the size of the torus through the three regimes of the model and its vacuum energy, and compares it with a 4-SPHERE of the same volume: curvature, levels, what a brane sees, Dirac zero modes, and how the size is held. Every equation used is written here. Section 25 of 7+1.nb contains 27 checks (27 True); the independent program tools/ref\_torus\_brane.py contains 37 (37 True); the 17 groups of numbers printed by both agree to a largest relative difference of $3.8\times10^{-10}$. The comparison with the sphere is checked in section 26 (33 checks, 33 True) and by tools/ref\_radion.py (31, 31 True).
\end{abstract}
\tableofcontents
\Needspace{12\baselineskip}\section*{1. Results at a glance}\addcontentsline{toc}{section}{1. Results at a glance}
{\footnotesize
\begin{longtable}{>{\raggedright\arraybackslash}p{\dimexpr 0.2200\linewidth-2\tabcolsep\relax}>{\raggedright\arraybackslash}p{\dimexpr 0.6600\linewidth-2\tabcolsep\relax}>{\raggedright\arraybackslash}p{\dimexpr 0.1200\linewidth-2\tabcolsep\relax}}
\toprule
\textbf{question} & \textbf{answer} & \textbf{where} \\
\midrule\endfirsthead
\toprule
\textbf{question} & \textbf{answer} & \textbf{where} \\
\midrule\endhead
\endfoot
\bottomrule\endlastfoot
What is the auxiliary torus? & the space of the four auxiliary coordinates $\vec y=(x_4,\dots,x_7)$ with $x_k\simeq x_k+L_cn_k$ ($n_k\in\mathbb Z$): $T^4=\mathbb R^4/(L_c\mathbb Z)^4$, compact, flat, without boundary & §3.1 \\
How big is it? & size $b$ with $b^8=\det(-g_{kl})$; each circle has the proper circumference $bL_c$; proper volume $V_\perp=b^4L_c^4$; only $bL_c$ is physical & §3.4, §3.5 \\
How big is it today? & $\ell_\perp=V_\perp^{1/4}=1560.45$ GeV$^{-1}=3.0792\times10^{-13}$ m $=307.9$ fm: about 366 proton radii, or the Bohr radius divided by 172 & §6.2 \\
What are its symmetries? & the translations $\vec y\to\vec y+\vec c$ (their charge is the Kaluza–Klein momentum) and the 384 signed permutations of the axes, which fix the world brane at $\vec y=0$ & §3.6 \\
What does a field on it look like? & a Fourier series in $\phi_{\vec n}=e^{i\vec q_{\vec n}\cdot\vec y}/\sqrt{V_\perp}$, $\vec q_{\vec n}=2\pi\vec n/L_c$, $\vec n\in\mathbb Z^4$ & §4.1 \\
What does ordinary spacetime see? & a Kaluza–Klein tower of masses $\sqrt{m^2+\mu_{\vec n}^2}$, $\mu_{\vec n}=2\pi|\vec n|/\ell_\perp$; the lightest non-zero one is 4.027 MeV; for dirac16complex four brane Dirac fermions per $\vec n$ & §4.4, §4.5 \\
How many levels? & $r_4(k)=8\sigma(k)-32\sigma(k/4)$ vectors with $|\vec n|^2=k$; $N(\mu)=V_\perp\mu^4/(32\pi^2)$ levels below $\mu$: $1.17\times10^{13}$ below 5 GeV, $1.88\times10^{26}$ below $M_8$ & §5 \\
What does it do to gravity? & $8\pi G_4=\kappa/V_\perp$; the force between two masses on the brane goes from $1/r^5$ ($r\ll\ell_\perp$) to $1/r$ ($r\gg\ell_\perp$); images $=$ Kaluza–Klein modes; $\frac53=1+\frac23$ (graviton $+$ radion) & §6 \\
How does its size evolve? & it contracts in the primordial regime, would re-expand ($b\propto\tau^{2/5}$) without the stabilizing potential, is caught by it, and stays frozen today and in the far future & §7 \\
How does the world brane see it? & the brane sits at the point $\vec y=0$; it breaks the translations, so it couples to every Kaluza–Klein level with the same strength $1/\sqrt{V_\perp}$ & §8 \\
What is its vacuum energy? & a Casimir energy, for a massless dirac16complex $6.8\times10^{34}\,\rho_\Lambda$: one more constant term of the cosmological constant & §9 \\
Why a torus and not a sphere? & a torus is flat: its size costs no curvature energy (a modulus until stabilized), it has Dirac zero modes, and a brane on it keeps 384 symmetries; a 4-sphere of the same volume (radius $1.359\times10^{-13}$ m) is curved, needs $(\Lambda+\kappa U_0)r_0^2=6$ to be held, has its first level at 2.903 MeV and no Dirac zero modes & §10 \\*
\end{longtable}}
\Needspace{12\baselineskip}\section*{2. Conventions}\addcontentsline{toc}{section}{2. Conventions}
Units $\hbar=c=1$ (energies, masses and inverse lengths in GeV; $\hbar c=1.973269804\times10^{-16}$ GeV m). The coordinates of 7+1 spacetime are $X=(x_0,x_1,\dots,x_7)$: $x_0$ is time (the direction of the octonion unit $e_0=1$), $x_1,x_2,x_3$ are ordinary space, and $x_4,\dots,x_7$ are the four auxiliary space-like coordinates, collected in $\vec y=(x_4,x_5,x_6,x_7)$ (components $y_k=x_k$, $k=4,\dots,7$); $\vec x=(x_1,x_2,x_3)$. Indices: $A,B=0,\dots,7$ for all eight directions, $\mu,\nu=0,\dots,3$ for the directions of the world brane, $k,l=4,\dots,7$ for the directions of the torus. A prime is $d/dx_0$. The frame metric is $\eta=\mathrm{diag}(+1,-1,\dots,-1)$, simply the squares of the octonion basis (explained in the next paragraph).

\textbf{The signature: $\{7,1,0\}$ and $\mathrm{diag}(+1,-1,\dots,-1)$.} The 7+1 spacetime has one time-like dimension ($x_0$, the direction of the octonion unit $e_0=1$) and seven space-like dimensions ($x_1,x_2,x_3$ ordinary space, $x_4,\dots,x_7$ auxiliary space). Two different objects describe this, and neither may be replaced by the other. (i) The tensor package xAct, used in section 12 of 7+1.nb exactly as in the author’s notebook Einstein-Lovelock-7+1-Nash.nb, declares the metric with \texttt{DefMetric[\{7, 1, 0\}, g71[-a, -b], CD, …]}: the list $\{7,1,0\}$ declares seven space-like dimensions, one time-like dimension and no degenerate dimension — the numbers of entries $+1$ and $-1$ of $\eta_{71}=\mathrm{diag}(-1,+1,\dots,+1)$ (\texttt{η71}), the same spacetime written in the sign convention of 4+4.nb (space-like $=+1$, time-like $=-1$). (ii) The matrix $\eta=-\eta_{71}=\mathrm{diag}(+1,-1,\dots,-1)$ (\texttt{η17}) is simply the squares of the octonion basis: $e_0^2=+1$ for the unit $e_0=1$ (time) and $e_i^2=-1$ for the seven imaginary units $e_1,\dots,e_7$ (space); more generally it is the real part of the product of two basis units, and it gives the quadratic form of the octonions:
\begin{gather*}
\eta_{AB}=\mathrm{Re}(e_Ae_B),\qquad e_0^2=+1,\qquad e_i^2=-1\ \ (i=1,\dots,7),\\ \mathrm{Re}(x\,x)=x_0^2-x_1^2-\cdots-x_7^2=\eta_{AB}\,x^Ax^B\qquad (x=x^Ae_A\in\mathbb O).
\end{gather*}
$\eta$ gives the signs of the Clifford relation $\Gamma_A\Gamma_B+\Gamma_B\Gamma_A=2\,\eta_{AB}\,1_{16}$ of the real 16×16 Dirac matrices: time squares to $+1_{16}$ and each of the seven space directions to $-1_{16}$, exactly like the octonion units (in the notation $Cl(p,q)$, with $p$ generators squaring to $+1$ and $q$ to $-1$, the algebra $Cl(1,7)\cong\mathbb R(16)$). It also gives the signs of the metric components: $g_{00}>0$ along the time-like direction, negative squares along the space-like directions. Every component computed by 7+1.nb follows from the component matrix of the metric; the declared list fixes for xAct the sign of the determinant of the metric, \texttt{SignDetOfMetric[g71]} $=-1$, which is the same in both sign conventions, since $\det \eta=\det \eta_{71}=-1$. 7+1.nb checks \texttt{SignatureOfMetric[g71]} $=\{7,1,0\}$, \texttt{SignDetOfMetric[g71]} $=-1$, $\mathrm{Re}(e_Ae_B)=\eta_{AB}$ and $\mathrm{Re}(x\,x)=\eta_{AB}x^Ax^B$.

\textbf{Physical constants} (as in sections 17–24 of 7+1.nb): the reduced eight-dimensional Planck mass $M_8=10$ TeV, with $\kappa=8\pi G_8=M_8^{-6}$ ($G_8$ the Newton constant of eight dimensions); the reduced four-dimensional Planck mass $M_{\rm Pl}=2.435\times10^{18}$ GeV, $8\pi G_4=M_{\rm Pl}^{-2}$; the radion mass $m_r=10$ GeV; the Higgs mass $m_h=125.20$ GeV.

\Needspace{12\baselineskip}\section*{3. The definition of the auxiliary torus}\addcontentsline{toc}{section}{3. The definition of the auxiliary torus}
\Needspace{8\baselineskip}\subsection*{3.1 Periodic coordinates, and the torus as a quotient space}\addcontentsline{toc}{subsection}{3.1 Periodic coordinates, and the torus as a quotient space}
\textbf{Definition.} The four auxiliary coordinates are periodic with the same coordinate period $L_c>0$:
\begin{equation*}
x_k\simeq x_k+L_c\,n_k\qquad(k=4,\dots,7,\quad n_k\in\mathbb Z).
\end{equation*}
Two points $\vec y$ and $\vec y\,'$ of $\mathbb R^4$ describe the SAME point of the auxiliary space when their difference is a vector of the lattice $\Lambda=L_c\mathbb Z^4$. The auxiliary space is the set of these classes — the quotient space
\begin{equation*}
T^4=\mathbb R^4/\Lambda,\qquad \Lambda=L_c\mathbb Z^4=\{L_c\vec n:\ \vec n\in\mathbb Z^4\},\qquad \vec y\sim\vec y\,'\iff\vec y-\vec y\,'\in\Lambda,
\end{equation*}
the four-dimensional torus. Every class has exactly one representative in the FUNDAMENTAL DOMAIN
\begin{equation*}
F=[0,L_c)^4\qquad(\text{equivalently }[-\tfrac{L_c}2,\tfrac{L_c}2)^4),
\end{equation*}
a hypercube whose opposite faces are glued together. The torus is the product of four circles, $T^4=S^1\times S^1\times S^1\times S^1$ with $S^1=\mathbb R/L_c\mathbb Z$. A function on the torus is the same thing as a function on $\mathbb R^4$ that is periodic in each coordinate,
\begin{equation*}
f(\vec y+L_c\vec e_k)=f(\vec y)\qquad(k=4,\dots,7),
\end{equation*}
($\vec e_k$ the unit vector along $y_k$). Three properties matter throughout: the torus is \textbf{compact} (it has a finite volume), it is \textbf{flat} (its metric below has zero curvature), and it has \textbf{no boundary} and \textbf{no special point} — every point looks like every other, because the translations $\vec y\to\vec y+\vec c$ are symmetries. That is why the world brane can be put at the point $\vec y=0$ without loss of generality.

\Needspace{8\baselineskip}\subsection*{3.2 Two examples to build intuition}\addcontentsline{toc}{subsection}{3.2 Two examples to build intuition}
\Needspace{6\baselineskip}\subsubsection*{Example 3.1 — the circle, and distances on the torus}
The circle $S^1=\mathbb R/L_c\mathbb Z$ is the interval $[0,L_c)$ with its two ends identified: walking past $y=L_c$ brings you back to $y=0$. With the metric $b^2dy^2$ (section 3.4) the proper distance between two points is the length of the SHORTER way around,
\begin{equation*}
d(y,y\,')=b\,\min_{n\in\mathbb Z}\,|y-y\,'-nL_c|,\qquad d(\vec y,\vec y\,')=b\,\min_{\vec n\in\mathbb Z^4}\,|\vec y-\vec y\,'-L_c\vec n|\ \ \text{on }T^4 .
\end{equation*}
For $y=0.9\,L_c$ and $y\,'=0.1\,L_c$ the distance is $b\cdot|0.9L_c-0.1L_c-L_c|=0.2\,bL_c$, not $0.8\,bL_c$; no two points of the circle are farther apart than $bL_c/2$, and no two points of $T^4$ farther than $\sqrt4\,bL_c/2=bL_c$.

\Needspace{6\baselineskip}\subsubsection*{Example 3.2 — the two-torus as a video-game screen}
The two-torus $T^2=\mathbb R^2/(L_c\mathbb Z)^2$ is a square whose left and right edges are glued, and whose top and bottom edges are glued: an object leaving through the right edge re-enters at the left edge at the same height, as in the old arcade screens. The square is flat — straight lines and the angles of triangles are those of the plane — although the surface closes on itself. The auxiliary torus is the same construction with four directions: a hypercube with its four pairs of opposite faces glued.

\Needspace{8\baselineskip}\subsection*{3.3 The spacetime of the model, and where the world brane is}\addcontentsline{toc}{subsection}{3.3 The spacetime of the model, and where the world brane is}
The spacetime of the model, the BULK, is the product
\begin{equation*}
M=\mathbb R_{x_0}\times\mathbb R^3_{\vec x}\times T^4_{\vec y}\qquad(\text{time}\times\text{ordinary space}\times\text{auxiliary torus}).
\end{equation*}
Gravity and the 16-component field dirac16complex live in the bulk. The Standard Model lives on the WORLD BRANE
\begin{equation*}
W=\{X\in M:\ x_4=x_5=x_6=x_7=0\}=\mathbb R_{x_0}\times\mathbb R^3_{\vec x}\times\{\vec y=0\},
\end{equation*}
the 3+1-dimensional submanifold through the point $\vec y=0$ of the torus. Its metric is the induced metric $h_{\mu\nu}(x)=g_{\mu\nu}(x,\vec y=0)$ ($\mu,\nu=0,\dots,3$). The brane is a point of the torus: everything said below about fields on the torus decides how the brane sees them (section 8).

\Needspace{8\baselineskip}\subsection*{3.4 The metric, the size \texorpdfstring{$b$}{b} and the proper volume \texorpdfstring{$V_\perp$}{V}}\addcontentsline{toc}{subsection}{3.4 The metric, the size \texorpdfstring{$b$}{b} and the proper volume \texorpdfstring{$V_\perp$}{V}}
Write the metric of the bulk as $ds^2=g_{AB}\,dX^AdX^B$ and its torus block as $g_{kl}$ ($k,l=4,\dots,7$; negative definite, since the four directions are space-like). \textbf{The size $b$ of the torus is defined by}
\begin{equation*}
b^8=\det(-g_{kl})\qquad(k,l=4,\dots,7),\qquad\text{so that}\qquad dV_T=\sqrt{\det(-g_{kl})}\,d^4y=b^4\,d^4y
\end{equation*}
is the proper volume element of the torus. All metrics of 7+1.nb that describe the cosmology of the model have the form
\begin{equation*}
ds^2=n^2dx_0^2-a^2\,(dx_1^2+dx_2^2+dx_3^2)-b^2\,(dx_4^2+dx_5^2+dx_6^2+dx_7^2),\qquad a=a(x_0),\ \ b=b(x_0),
\end{equation*}
($n$ the lapse, $n=1$ in the cosmic time $x_0$; $a$ the scale factor of ordinary space), so $-g_{kl}=b^2\delta_{kl}$, and $b$ is the scale factor of the auxiliary space. Then each of the four circles has the proper circumference, and the torus the proper volume,
\begin{equation*}
C=\int_0^{L_c}b\,dx_k=b\,L_c,\qquad V_\perp=\int_F b^4\,d^4y=b^4L_c^4,\qquad \ell_\perp=V_\perp^{1/4}=b\,L_c .
\end{equation*}
Today the size is frozen at $b=b_0$, so $V_\perp=b_0^4L_c^4$ and $\ell_\perp=b_0L_c$.

\Needspace{6\baselineskip}\subsubsection*{Example 3.3 — the volume of the torus for the cosmological metric}
For the metric above the torus block is $g_{kl}=-b^2\delta_{kl}$, so $\det(-g_{kl})=b^8$, and
\begin{equation*}
V_\perp=\int_0^{L_c}\!\!\int_0^{L_c}\!\!\int_0^{L_c}\!\!\int_0^{L_c}\sqrt{\det(-g_{kl})}\;dx_4\,dx_5\,dx_6\,dx_7=b^4L_c^4,
\end{equation*}
as checked symbolically in section 25.1 of 7+1.nb (and by the Python program).

\Needspace{8\baselineskip}\subsection*{3.5 Only the product \texorpdfstring{$bL_c$}{bLc} is physical}\addcontentsline{toc}{subsection}{3.5 Only the product \texorpdfstring{$bL_c$}{bLc} is physical}
The coordinate $\vec y$ is a label; a rescaled label describes the same torus. Under $\vec y=\lambda\vec y\,'$ ($\lambda>0$):
\begin{equation*}
b^2\,d\vec y\cdot d\vec y=(\lambda b)^2\,d\vec y\,'\cdot d\vec y\,',\qquad L_c'=\frac{L_c}\lambda,\qquad b'=\lambda b,\qquad b'L_c'=bL_c,\qquad b'^4L_c'^4=b^4L_c^4 .
\end{equation*}
So the period $L_c$ and the size $b$ are not separately measurable — only the proper circumference $bL_c$ and the volume $V_\perp$ are. A common choice is $b_0=1$, which makes $L_c=\ell_\perp$ the proper size today; 7+1.nb keeps both symbols. \textbf{Example 3.4:} the torus with $L_c=1$ (in any unit of coordinate length) and $b_0=1560.45$ GeV$^{-1}$ per unit, and the torus with $L_c=1560.45$ GeV$^{-1}$ and $b_0=1$, are the same torus (checked symbolically in section 25.1).

\Needspace{8\baselineskip}\subsection*{3.6 Symmetries}\addcontentsline{toc}{subsection}{3.6 Symmetries}
The isometries of the flat torus (the maps that preserve all distances) are of two kinds:

\begin{itemize}
\item the \textbf{translations} $\vec y\to\vec y+\vec c$ ($\vec c\in\mathbb R^4$, modulo $\Lambda$). They form the group $U(1)^4$; the conserved charge that belongs to them is the momentum along the torus, the Kaluza–Klein momentum $\vec n$ (section 8);
\item the \textbf{signed permutations} of the four axes, $\vec y\to O\vec y$ with $O$ a $4\times4$ matrix that has one entry $\pm1$ in each row and column. These are exactly the orthogonal matrices that map the lattice $\Lambda$ onto itself (an orthogonal matrix with integer entries is a signed permutation). There are $2^4\cdot4!=384$ of them (the hyperoctahedral group $B_4$); they fix the point $\vec y=0$ and with it the world brane.
\end{itemize}
A rotation by a general angle is NOT a symmetry of the torus: it does not map the lattice to itself. The torus is flat and locally looks like $\mathbb R^4$, but globally it is not isotropic. The signed permutations permute the vectors $\vec n\in\mathbb Z^4$ of each shell $|\vec n|^2=k$, and two vectors are in the same orbit exactly when they have the same multiset of $|n_i|$.

\Needspace{6\baselineskip}\subsubsection*{Example 3.5 — the shells of the lattice and their orbits}
{\footnotesize
\begin{longtable}{>{\raggedright\arraybackslash}p{\dimexpr 0.1400\linewidth-2\tabcolsep\relax}>{\raggedright\arraybackslash}p{\dimexpr 0.1200\linewidth-2\tabcolsep\relax}>{\raggedright\arraybackslash}p{\dimexpr 0.2000\linewidth-2\tabcolsep\relax}>{\raggedright\arraybackslash}p{\dimexpr 0.5400\linewidth-2\tabcolsep\relax}}
\toprule
\textbf{$k=|\vec n|^2$} & \textbf{$r_4(k)$} & \textbf{orbit sizes} & \textbf{one vector of each orbit (in the same order)} \\
\midrule\endfirsthead
\toprule
\textbf{$k=|\vec n|^2$} & \textbf{$r_4(k)$} & \textbf{orbit sizes} & \textbf{one vector of each orbit (in the same order)} \\
\midrule\endhead
\endfoot
\bottomrule\endlastfoot
1 & 8 & 8 & $(1,0,0,0)$ \\
2 & 24 & 24 & $(1,1,0,0)$ \\
3 & 32 & 32 & $(1,1,1,0)$ \\
4 & 24 & 8 + 16 & $(2,0,0,0)$, $(1,1,1,1)$ \\
5 & 48 & 48 & $(2,1,0,0)$ \\
6 & 96 & 96 & $(2,1,1,0)$ \\
7 & 64 & 64 & $(2,1,1,1)$ \\
8 & 24 & 24 & $(2,2,0,0)$ \\
9 & 104 & 8 + 96 & $(3,0,0,0)$, $(2,2,1,0)$ \\
10 & 144 & 48 + 96 & $(3,1,0,0)$, $(2,2,1,1)$ \\
11 & 96 & 96 & $(3,1,1,0)$ \\
12 & 96 & 32 + 64 & $(2,2,2,0)$, $(3,1,1,1)$ \\*
\end{longtable}}
For example the shell $k=4$ holds the 8 vectors $\pm2\vec e_i$ and the 16 vectors $(\pm1,\pm1,\pm1,\pm1)$: they have the same length but are not related by a symmetry. The orbit sizes are printed by section 25.1 of 7+1.nb and agree with the Python program.

\Needspace{12\baselineskip}\section*{4. Fields on the torus: Fourier modes and the Kaluza–Klein tower}\addcontentsline{toc}{section}{4. Fields on the torus: Fourier modes and the Kaluza–Klein tower}
\Needspace{8\baselineskip}\subsection*{4.1 Fourier modes}\addcontentsline{toc}{subsection}{4.1 Fourier modes}
A field on the torus is periodic in each $y_k$. A plane wave $e^{i\vec q\cdot\vec y}$ is periodic only if $e^{iq_kL_c}=1$, i.e. $q_k=2\pi n_k/L_c$ with $n_k\in\mathbb Z$: the allowed wave vectors form the DUAL LATTICE $(2\pi/L_c)\mathbb Z^4$. The normalized MODES are
\begin{equation*}
\phi_{\vec n}(\vec y)=\frac{e^{i\vec q_{\vec n}\cdot\vec y}}{\sqrt{V_\perp}},\qquad \vec q_{\vec n}=\frac{2\pi\vec n}{L_c},\qquad \vec n\in\mathbb Z^4 .
\end{equation*}
($\vec q_{\vec n}$ is the wave vector in the coordinate $\vec y$; the physical momentum along the torus is $\vec q_{\vec n}/b$, because the proper distance is $b\,|d\vec y|$.) They are ORTHONORMAL for the proper volume element and COMPLETE:
\begin{align*}
\int_F b^4d^4y\;\phi_{\vec n}^*\,\phi_{\vec m}&=\delta_{\vec n\vec m}, \\
f(\vec y)=\sum_{\vec n\in\mathbb Z^4}c_{\vec n}\,\phi_{\vec n}(\vec y),\qquad c_{\vec n}&=\int_F b^4d^4y\;\phi_{\vec n}^*f,\qquad \int_F b^4d^4y\;|f|^2=\sum_{\vec n}|c_{\vec n}|^2\quad(\text{Parseval}).
\end{align*}
The proof of orthonormality is one line per direction: $\int_0^{L_c}e^{2\pi i(m_k-n_k)y_k/L_c}dy_k=L_c\,\delta_{m_kn_k}$, and $b^4L_c^4/V_\perp=1$.

\Needspace{8\baselineskip}\subsection*{4.2 Example: orthonormality and Parseval’s identity}\addcontentsline{toc}{subsection}{4.2 Example: orthonormality and Parseval’s identity}
\Needspace{6\baselineskip}\subsubsection*{Example 4.1 — three exact integrals}
Section 25.2 of 7+1.nb integrates exactly over the fundamental domain: $\int_Fb^4\phi^*_{\vec n}\phi_{\vec m}=1$ for $\vec n=\vec m=(1,0,-2,3)$, $=0$ for $\vec n=(1,0,-2,3)$, $\vec m=(1,1,-2,3)$, and $=0$ for $\vec n=0$, $\vec m=(0,2,0,0)$.

\Needspace{6\baselineskip}\subsubsection*{Example 4.2 — a smooth bump and its 81 Fourier coefficients}
Take $f(\vec y)=\prod_{k=4}^7\cos^2(\pi y_k/L_c)$, a bump that is largest on the brane ($f=1$ at $\vec y=0$) and vanishes on the faces $y_k=\pm L_c/2$. Since $\cos^2\theta=\frac12+\frac14e^{2i\theta}+\frac14e^{-2i\theta}$,
\begin{equation*}
f=\prod_{k=4}^7\Big(\frac12+\frac14e^{2\pi iy_k/L_c}+\frac14e^{-2\pi iy_k/L_c}\Big),\qquad c_{\vec n}=\sqrt{V_\perp}\prod_{k}w(n_k),\quad w(0)=\tfrac12,\ w(\pm1)=\tfrac14,
\end{equation*}
so exactly $3^4=81$ coefficients are non-zero ($n_k\in\{-1,0,1\}$). Parseval’s identity holds:
\begin{equation*}
\sum_{\vec n}|c_{\vec n}|^2=V_\perp\Big(\frac14+\frac1{16}+\frac1{16}\Big)^4=\Big(\frac38\Big)^4V_\perp=\int_Fb^4d^4y\,f^2,
\end{equation*}
because $\int_0^{L_c}\cos^4(\pi y/L_c)\,dy=\frac38L_c$. Section 25.2 checks three coefficients by integration and both sides of the identity.

\Needspace{8\baselineskip}\subsection*{4.3 The Laplacian of the torus}\addcontentsline{toc}{subsection}{4.3 The Laplacian of the torus}
The Laplacian of the flat torus with the metric $b^2\delta_{kl}$ is
\begin{equation*}
\Delta_T=\frac1{b^2}\sum_{k=4}^7\frac{\partial^2}{\partial y_k^2},\qquad -\Delta_T\,\phi_{\vec n}=\frac{|\vec q_{\vec n}|^2}{b^2}\,\phi_{\vec n}=\Big(\frac{2\pi}{bL_c}\Big)^2|\vec n|^2\,\phi_{\vec n} :
\end{equation*}
the modes are its eigenfunctions, and the eigenvalue depends only on $k=|\vec n|^2$ (with the multiplicity $r_4(k)$ of section 5). \textbf{Example 4.3:} for $\vec n=(2,-1,0,3)$, $|\vec n|^2=14$ and $-\Delta_T\phi_{\vec n}=(2\pi/(bL_c))^2\cdot14\,\phi_{\vec n}$ (section 25.2).

\Needspace{8\baselineskip}\subsection*{4.4 Kaluza–Klein masses: a scalar field}\addcontentsline{toc}{subsection}{4.4 Kaluza–Klein masses: a scalar field}
Take a static torus ($b=b_0$) over flat ordinary spacetime, $g=\mathrm{diag}(1,-1,-1,-1,-b_0^2,-b_0^2,-b_0^2,-b_0^2)$, and a scalar field $\Phi$ of mass $m$ in the bulk, with the Klein–Gordon equation
\begin{equation*}
(\Box+m^2)\Phi=0,\qquad \Box\Phi=\frac1{\sqrt{|g|}}\,\partial_A\big(\sqrt{|g|}\,g^{AB}\partial_B\Phi\big).
\end{equation*}
For $\Phi=e^{-iEx_0+i\vec k\cdot\vec x}\,\phi_{\vec n}(\vec y)$ the operator gives a number times $\Phi$:
\begin{align*}
(\Box+m^2)\Phi&=\Big(-E^2+|\vec k|^2+\frac{|\vec q_{\vec n}|^2}{b_0^2}+m^2\Big)\Phi=0\quad\Longrightarrow\quad E^2=|\vec k|^2+M_{\vec n}^2, \\
M_{\vec n}&=\sqrt{m^2+\mu_{\vec n}^2},\qquad \mu_{\vec n}=\frac{|\vec q_{\vec n}|}{b_0}=\frac{2\pi|\vec n|}{\ell_\perp}.
\end{align*}
An observer who cannot resolve the torus sees, for every $\vec n$, a particle of mass $M_{\vec n}$: \textbf{momentum along the hidden directions looks like mass}. The set of these masses is the KALUZA–KLEIN TOWER; $\vec n=0$ is the ZERO MODE (mass $m$, constant over the torus). Section 25.2 checks the dispersion relation symbolically for $\vec n=(1,-2,0,2)$.

\Needspace{8\baselineskip}\subsection*{4.5 Kaluza–Klein masses: dirac16complex}\addcontentsline{toc}{subsection}{4.5 Kaluza–Klein masses: dirac16complex}
For the 16-component field the same happens with four 4-component spinors per mode. With the $4\times4$ Dirac matrices (Dirac representation) and four anti-Hermitian matrices $\rho_k$,
\begin{align*}
\gamma_0&=\mathrm{diag}(1,1,-1,-1),\quad \gamma_i=\begin{pmatrix}0&\sigma_i\\-\sigma_i&0\end{pmatrix}\ (i=1,2,3),\quad \gamma_5=i\gamma_0\gamma_1\gamma_2\gamma_3,\quad \gamma_\mu\gamma_\nu+\gamma_\nu\gamma_\mu=2\eta_{\mu\nu}1_4, \\
\rho_4&=i\gamma_0,\quad \rho_5=-\gamma_1,\quad \rho_6=-\gamma_2,\quad \rho_7=-\gamma_3,\qquad \rho_k\rho_l+\rho_l\rho_k=-2\delta_{kl}1_4,
\end{align*}
($\sigma_i$ the Pauli matrices), the sixteen matrices $R_\mu=\gamma_\mu\otimes1_4$ ($\mu=0,\dots,3$) and $R_k=\gamma_5\otimes\rho_k$ ($k=4,\dots,7$) satisfy the same Clifford relation $R_AR_B+R_BR_A=2\eta_{AB}1_{16}$ as the real $16\times16$ Dirac matrices $\Gamma_A$ of dirac16complex, so there is a constant unitary $16\times16$ matrix $U$ with $\Gamma_AU=UR_A$ for all $A$ (both representations are irreducible). Writing the field as $\psi=U\zeta$, $\zeta\in\mathbb C^4\otimes\mathbb C^4$, separates a 4-component spinor index of the brane from a 4-valued index $a$ of the transverse directions. For $\psi\propto e^{-iEx_0+i\vec p\cdot\vec x+i\vec q\cdot\vec y}$ at $b=b_0$ the Dirac equation $(i\Gamma^Ae_A{}^B\partial_B-m)\psi=0$ of dirac16complex ($\Gamma^A=\eta^{AA}\Gamma_A$, frame $e_A{}^B=\mathrm{diag}(1,1,1,1,b_0^{-1},b_0^{-1},b_0^{-1},b_0^{-1})$) becomes
\begin{equation*}
\big(\gamma\cdot p\otimes1_4+X-m\big)\zeta=0,\qquad \gamma\cdot p=E\gamma_0-\sum_{i=1}^3p_i\gamma_i,\qquad X=\gamma_5\otimes\frac{\vec\rho\cdot\vec q}{b_0},\qquad X^2=-\frac{|\vec q|^2}{b_0^2}\,1_{16},
\end{equation*}
and $X$ anticommutes with every $\gamma_\mu\otimes1_4$. Squaring, $E^2=|\vec p|^2+m^2+|\vec q|^2/b_0^2$: each momentum vector $\vec n$ of the torus gives FOUR Dirac fermions of mass $M_{\vec n}$ in ordinary spacetime (section 18.1 of 7+1.nb).

\Needspace{6\baselineskip}\subsubsection*{Example 4.4 — the lowest levels today}
{\footnotesize
\begin{longtable}{>{\raggedright\arraybackslash}p{\dimexpr 0.1600\linewidth-2\tabcolsep\relax}>{\raggedright\arraybackslash}p{\dimexpr 0.2600\linewidth-2\tabcolsep\relax}>{\raggedright\arraybackslash}p{\dimexpr 0.2600\linewidth-2\tabcolsep\relax}>{\raggedright\arraybackslash}p{\dimexpr 0.3200\linewidth-2\tabcolsep\relax}}
\toprule
\textbf{$k=|\vec n|^2$} & \textbf{momentum vectors $r_4(k)$} & \textbf{Dirac fermions $4r_4(k)$} & \textbf{mass $2\pi\sqrt k/\ell_\perp$ (MeV), $m=0$} \\
\midrule\endfirsthead
\toprule
\textbf{$k=|\vec n|^2$} & \textbf{momentum vectors $r_4(k)$} & \textbf{Dirac fermions $4r_4(k)$} & \textbf{mass $2\pi\sqrt k/\ell_\perp$ (MeV), $m=0$} \\
\midrule\endhead
\endfoot
\bottomrule\endlastfoot
1 & 8 & 32 & 4.02652 \\
2 & 24 & 96 & 5.69437 \\
3 & 32 & 128 & 6.97415 \\
4 & 24 & 96 & 8.05305 \\
5 & 48 & 192 & 9.00358 \\
6 & 96 & 384 & 9.86293 \\
7 & 64 & 256 & 10.6532 \\
8 & 24 & 96 & 11.3887 \\*
\end{longtable}}
With $m=0$ there are 4 massless Dirac fermions (the zero mode), then 32 of mass 4.0265 MeV, 96 of mass 5.6944 MeV, and so on; the spacing of the masses shrinks like $1/\sqrt k$, and the number of states grows like $k$ (section 5).

\Needspace{8\baselineskip}\subsection*{4.6 Fermions on the torus: the boundary conditions}\addcontentsline{toc}{subsection}{4.6 Fermions on the torus: the boundary conditions}
A spinor need not return to itself when it goes once around a circle: it may return to itself or change sign. 7+1.nb uses PERIODIC conditions on all four circles (the modes $e^{i\vec q_{\vec n}\cdot\vec y}$ above; also in the Casimir energy of section 9). It is one of the $2^4=16$ possible choices (spin structures of $T^4$); with an antiperiodic circle $k$ the momenta along it would be shifted to $q_k=2\pi(n_k+\frac12)/L_c$ and there would be no zero mode.

\Needspace{8\baselineskip}\subsection*{4.7 The graviton on the torus}\addcontentsline{toc}{subsection}{4.7 The graviton on the torus}
Gravity is also a bulk field. A massless graviton in $D$ spacetime dimensions has $D(D-3)/2$ states: 20 for $D=8$. Seen from ordinary spacetime they split as follows (counting the states of the little groups):
\begin{align*}
\text{zero mode (massless): }&20=\underbrace{2}_{\text{graviton }g_{\mu\nu}}+\underbrace{4\times2}_{\text{4 vectors }g_{\mu k}}+\underbrace{10}_{\text{scalars }g_{kl}}, \\
\text{each massive level: }&20=\underbrace{5}_{\text{massive spin 2}}+\underbrace{3\times3}_{\text{3 massive vectors}}+\underbrace{6}_{\text{scalars}}.
\end{align*}
The ten massless scalars $g_{kl}$ are the MODULI of the torus: one is its volume (the RADION, section 6.5), nine describe its shape (the angles and ratios of the circles). The four vectors $g_{\mu k}$ are the graviphotons. The stabilizing potential $U(b)$ of 7+1.nb depends on the volume only; the shape moduli and the graviphotons are not given a potential or studied in 7+1.nb. Fields on the brane feel only the induced metric $h_{\mu\nu}=g_{\mu\nu}(x,0)$, which contains neither of them (section 11).

\Needspace{12\baselineskip}\section*{5. Counting the levels}\addcontentsline{toc}{section}{5. Counting the levels}
\Needspace{8\baselineskip}\subsection*{5.1 Jacobi’s four-square theorem}\addcontentsline{toc}{subsection}{5.1 Jacobi’s four-square theorem}
The masses depend only on $k=|\vec n|^2$. The number of $\vec n\in\mathbb Z^4$ with $|\vec n|^2=k$ is given by Jacobi’s four-square theorem (1829):
\begin{equation*}
r_4(k)=\#\{\vec n\in\mathbb Z^4:|\vec n|^2=k\}=8\sigma(k)-32\sigma(k/4)\quad(k\ge1),\qquad r_4(0)=1,\qquad \sigma(k)=\sum_{d\mid k}d,
\end{equation*}
with $\sigma(k/4)=0$ unless 4 divides $k$. \textbf{Example 5.1:} $r_4(4)=8\,\sigma(4)-32\,\sigma(1)=8\cdot7-32\cdot1=24=8+16$ (the two orbits of Example 3.5); $r_4(9)=8\cdot13=104=8+96$. The first values:

{\footnotesize
\begin{longtable}{>{\raggedright\arraybackslash}p{\dimexpr 0.1000\linewidth-2\tabcolsep\relax}>{\raggedright\arraybackslash}p{\dimexpr 0.0529\linewidth-2\tabcolsep\relax}>{\raggedright\arraybackslash}p{\dimexpr 0.0529\linewidth-2\tabcolsep\relax}>{\raggedright\arraybackslash}p{\dimexpr 0.0529\linewidth-2\tabcolsep\relax}>{\raggedright\arraybackslash}p{\dimexpr 0.0529\linewidth-2\tabcolsep\relax}>{\raggedright\arraybackslash}p{\dimexpr 0.0529\linewidth-2\tabcolsep\relax}>{\raggedright\arraybackslash}p{\dimexpr 0.0529\linewidth-2\tabcolsep\relax}>{\raggedright\arraybackslash}p{\dimexpr 0.0529\linewidth-2\tabcolsep\relax}>{\raggedright\arraybackslash}p{\dimexpr 0.0529\linewidth-2\tabcolsep\relax}>{\raggedright\arraybackslash}p{\dimexpr 0.0529\linewidth-2\tabcolsep\relax}>{\raggedright\arraybackslash}p{\dimexpr 0.0529\linewidth-2\tabcolsep\relax}>{\raggedright\arraybackslash}p{\dimexpr 0.0529\linewidth-2\tabcolsep\relax}>{\raggedright\arraybackslash}p{\dimexpr 0.0529\linewidth-2\tabcolsep\relax}>{\raggedright\arraybackslash}p{\dimexpr 0.0529\linewidth-2\tabcolsep\relax}>{\raggedright\arraybackslash}p{\dimexpr 0.0529\linewidth-2\tabcolsep\relax}>{\raggedright\arraybackslash}p{\dimexpr 0.0529\linewidth-2\tabcolsep\relax}>{\raggedright\arraybackslash}p{\dimexpr 0.0529\linewidth-2\tabcolsep\relax}>{\raggedright\arraybackslash}p{\dimexpr 0.0529\linewidth-2\tabcolsep\relax}}
\toprule
\textbf{$k$} & \textbf{0} & \textbf{1} & \textbf{2} & \textbf{3} & \textbf{4} & \textbf{5} & \textbf{6} & \textbf{7} & \textbf{8} & \textbf{9} & \textbf{10} & \textbf{11} & \textbf{12} & \textbf{13} & \textbf{14} & \textbf{15} & \textbf{16} \\
\midrule\endfirsthead
\toprule
\textbf{$k$} & \textbf{0} & \textbf{1} & \textbf{2} & \textbf{3} & \textbf{4} & \textbf{5} & \textbf{6} & \textbf{7} & \textbf{8} & \textbf{9} & \textbf{10} & \textbf{11} & \textbf{12} & \textbf{13} & \textbf{14} & \textbf{15} & \textbf{16} \\
\midrule\endhead
\endfoot
\bottomrule\endlastfoot
$r_4(k)$ & 1 & 8 & 24 & 32 & 24 & 48 & 96 & 64 & 24 & 104 & 144 & 96 & 96 & 112 & 192 & 192 & 24 \\*
\end{longtable}}
Section 25.3 counts all $\vec n\in\{-10,\dots,10\}^4$ directly and finds $r_4(k)$ for every $k\le100$.

\Needspace{8\baselineskip}\subsection*{5.2 Weyl’s law: the number of levels grows like the volume of a ball}\addcontentsline{toc}{subsection}{5.2 Weyl’s law: the number of levels grows like the volume of a ball}
The number of lattice vectors in a ball of radius $R$ approaches the volume $\frac{\pi^2}2R^4$ of the four-dimensional ball:
\begin{equation*}
N(R)=\#\{\vec n\in\mathbb Z^4:|\vec n|\le R\}=\sum_{k=0}^{R^2}r_4(k)\;\simeq\;\frac{\pi^2}2R^4\qquad(R\to\infty).
\end{equation*}
{\footnotesize
\begin{longtable}{>{\raggedright\arraybackslash}p{\dimexpr 0.1200\linewidth-2\tabcolsep\relax}>{\raggedright\arraybackslash}p{\dimexpr 0.3000\linewidth-2\tabcolsep\relax}>{\raggedright\arraybackslash}p{\dimexpr 0.3000\linewidth-2\tabcolsep\relax}>{\raggedright\arraybackslash}p{\dimexpr 0.2800\linewidth-2\tabcolsep\relax}}
\toprule
\textbf{$R$} & \textbf{$N(R)$ (lattice count)} & \textbf{$\pi^2R^4/2$} & \textbf{$N(R)/(\pi^2R^4/2)-1$} \\
\midrule\endfirsthead
\toprule
\textbf{$R$} & \textbf{$N(R)$ (lattice count)} & \textbf{$\pi^2R^4/2$} & \textbf{$N(R)/(\pi^2R^4/2)-1$} \\
\midrule\endhead
\endfoot
\bottomrule\endlastfoot
10 & 49689 & 49348.02201 & 0.00691 \\
30 & 4000425 & $3.997189782\times10^{6}$ & 0.000809 \\
100 & 493490641 & $4.934802201\times10^{8}$ & $2.11\times10^{-5}$ \\
300 & 39972070929 & $3.997189782\times10^{10}$ & $4.33\times10^{-6}$ \\
1000 & 4934805110729 & $4.934802201\times10^{12}$ & $5.90\times10^{-7}$ \\*
\end{longtable}}
The relative deviation falls roughly like $1/R^2$ (the lattice points near the surface of the ball): Weyl’s law (1911) for the eigenvalues of the Laplacian.

\Needspace{8\baselineskip}\subsection*{5.3 Levels below a mass, and sums over levels}\addcontentsline{toc}{subsection}{5.3 Levels below a mass, and sums over levels}
A level with $\mu_{\vec n}\le\mu$ has $|\vec n|\le R=\mu\ell_\perp/(2\pi)$, so the number of levels below $\mu$ and their density are
\begin{equation*}
N(\mu)\simeq\frac{\pi^2}2\Big(\frac{\mu\ell_\perp}{2\pi}\Big)^4=\frac{V_\perp\mu^4}{32\pi^2},\qquad dN=\frac{V_\perp\mu^3\,d\mu}{8\pi^2},\qquad \sum_{\vec n}F(\mu_{\vec n})\simeq\frac{V_\perp}{8\pi^2}\int_0^\infty\mu^3F(\mu)\,d\mu .
\end{equation*}
\Needspace{6\baselineskip}\subsubsection*{Example 5.2 — how many levels?}
With $V_\perp=5.929\times10^{12}$ GeV$^{-4}$ (section 6.2): $N(5\ \text{GeV})=1.17335\times10^{13}$ (the lattice count, section 22.2 of 7+1.nb, is 11733511217161: they agree to $5\times10^{-7}$); $N(m_h)=4.61\times10^{18}$; $N(M_8)=M_{\rm Pl}^2/(32\pi^2M_8^2)=1.88\times10^{26}$, since $V_\perp M_8^4=M_{\rm Pl}^2/M_8^2$.

\Needspace{6\baselineskip}\subsubsection*{Example 5.3 — a weighted sum: the levels open to the decay of a Higgs boson}
A Higgs boson can decay into a neutrino and any level lighter than $m_h$, with a phase-space weight $(1-\mu_{\vec n}^2/m_h^2)^2$ (section 24.4 of 7+1.nb). The weighted number of levels is
\begin{equation*}
\sum_{\mu_{\vec n}<m_h}\Big(1-\frac{\mu_{\vec n}^2}{m_h^2}\Big)^2\simeq\frac{V_\perp}{8\pi^2}\int_0^{m_h}\mu^3\Big(1-\frac{\mu^2}{m_h^2}\Big)^2d\mu=\frac{V_\perp m_h^4}{8\pi^2}\int_0^1u^3(1-u^2)^2du=\frac{V_\perp m_h^4}{192\pi^2}=7.688\times10^{17},
\end{equation*}
since $\int_0^1u^3(1-u^2)^2du=\frac14-\frac13+\frac18=\frac1{24}$. It is smaller than the number $N(m_h)=4.61\times10^{18}$ of levels below $m_h$ because the heavy levels have little phase space. (For a lattice radius $R=100$ the lattice sum agrees with the integral to $10^{-8}$; section 24.4.)

\Needspace{12\baselineskip}\section*{6. The size of the torus today, and four-dimensional gravity}\addcontentsline{toc}{section}{6. The size of the torus today, and four-dimensional gravity}
\Needspace{8\baselineskip}\subsection*{6.1 From eight to four dimensions: \texorpdfstring{$8\pi G_4=\kappa/V_\perp$}{8 G4=/V}}\addcontentsline{toc}{subsection}{6.1 From eight to four dimensions: \texorpdfstring{$8\pi G_4=\kappa/V_\perp$}{8 G4=/V}}
The Einstein–Hilbert term of the bulk is $S_1=-\frac1{2\kappa}\int d^8X\sqrt{|g|}\,R$ ($\kappa=8\pi G_8$; the curvature conventions of 7+1.nb). For a static torus over a four-dimensional spacetime, $g=\mathrm{diag}\big(g^{(4)}_{\mu\nu}(x),\,-b_0^2\delta_{kl}\big)$, the volume element is $\sqrt{|g|}=b_0^4\sqrt{|g^{(4)}|}$ and, the torus being flat and static, the curvature of the product is that of the four-dimensional factor, $R=R^{(4)}$. Integrating over the torus,
\begin{align*}
S_1&=-\frac{b_0^4L_c^4}{2\kappa}\int d^4x\sqrt{|g^{(4)}|}\,R^{(4)}=-\frac1{16\pi G_4}\int d^4x\sqrt{|g^{(4)}|}\,R^{(4)}, \\
8\pi G_4&=\frac{\kappa}{V_\perp},\qquad M_{\rm Pl}^2=\frac1{8\pi G_4}=\frac{V_\perp}\kappa=V_\perp M_8^6 .
\end{align*}
\textbf{Example 6.1:} for the cosmological metric 7+1.nb finds $R=-2(3H_a'+6H_a^2+4H_b'+10H_b^2+12H_aH_b)$ with $H_a=a'/a$, $H_b=b'/b$ (section 11.2 of 7+1.nb); for a static torus ($H_b=0$) this is $-6(H_a'+2H_a^2)$, the scalar curvature of the four-dimensional Friedmann–Robertson–Walker metric $\mathrm{diag}(1,-a^2,-a^2,-a^2)$ in the same conventions (checked by the Python program), and $\sqrt{|g|}=a^3b^4$. A bigger torus means weaker four-dimensional gravity: $G_4=G_8/V_\perp$.

\Needspace{8\baselineskip}\subsection*{6.2 The numbers today}\addcontentsline{toc}{subsection}{6.2 The numbers today}
With $M_8=10$ TeV and $M_{\rm Pl}=2.435\times10^{18}$ GeV, $V_\perp=M_{\rm Pl}^2/M_8^6$. For comparison (CODATA 2018): the proton charge radius $r_p=0.8414$ fm and the Bohr radius $a_0=5.29177210903\times10^{-11}$ m.

{\footnotesize
\begin{longtable}{>{\raggedright\arraybackslash}p{\dimexpr 0.3600\linewidth-2\tabcolsep\relax}>{\raggedright\arraybackslash}p{\dimexpr 0.3400\linewidth-2\tabcolsep\relax}>{\raggedright\arraybackslash}p{\dimexpr 0.3000\linewidth-2\tabcolsep\relax}}
\toprule
\textbf{quantity} & \textbf{value} & \textbf{formula} \\
\midrule\endfirsthead
\toprule
\textbf{quantity} & \textbf{value} & \textbf{formula} \\
\midrule\endhead
\endfoot
\bottomrule\endlastfoot
proper volume & $V_\perp=5.929225\times10^{12}$ GeV$^{-4}$ & $M_{\rm Pl}^2/M_8^6$ \\
proper size (circumference of each circle) & $\ell_\perp=1560.4487$ GeV$^{-1}$ & $V_\perp^{1/4}=b_0L_c$ \\
the same in metres & $3.07919\times10^{-13}$ m $=307.919$ fm & $\ell_\perp\cdot\hbar c$ \\
compared with the proton & $\ell_\perp/r_p=365.96$ &  \\
compared with the hydrogen atom & $a_0/\ell_\perp=171.86$ &  \\
lowest Kaluza–Klein mass & $2\pi/\ell_\perp=4.02652$ MeV & $\mu_{\vec n}$ for $|\vec n|=1$ \\
radius of each circle (convention of the Particle Data Group, $V_\perp=(2\pi R_{\rm PDG})^4$) & $R_{\rm PDG}=248.353$ GeV$^{-1}=4.9007\times10^{-14}$ m & $\ell_\perp/(2\pi)$ \\
fundamental scale in that convention & $M_D=34.050$ TeV & $M_{\rm Pl}^2=R_{\rm PDG}^4M_D^6\Rightarrow M_D=(2\pi)^{2/3}M_8$ \\*
\end{longtable}}
So the auxiliary torus is about 366 times larger than a proton and 172 times smaller than a hydrogen atom. Searches for large extra dimensions at colliders exclude $M_D$ below 5.9 to 11.2 TeV for 6 to 2 extra dimensions (Particle Data Group 2024); $M_D=34.05$ TeV is above these bounds. The same numbers are printed by sections 17.5 and 22.2 of 7+1.nb.

\Needspace{8\baselineskip}\subsection*{6.3 Newton’s law on the torus: images and Kaluza–Klein modes (worked example)}\addcontentsline{toc}{subsection}{6.3 Newton’s law on the torus: images and Kaluza–Klein modes (worked example)}
This example shows, in one formula, how a compact dimension hides itself. Work in LINEARIZED eight-dimensional gravity WITHOUT the stabilizing potential (section 6.4 adds it).

\textbf{Step 1: Newton’s law in eight dimensions.} In $D$ spacetime dimensions the Newtonian potential of a static mass density $\rho$ obeys (the potential $\Phi_{\rm N}$ of a static weak field is defined by $g_{00}=1+2\Phi_{\rm N}$; this follows from the linearized Einstein equation)
\begin{equation*}
\nabla^2\Phi_{\rm N}=8\pi G_D\,\frac{D-3}{D-2}\,\rho,\qquad \nabla^2\frac1{r^{d-2}}=-(d-2)\,\Omega_{d-1}\,\delta^d(\vec r),\qquad \Omega_{d-1}=\frac{2\pi^{d/2}}{\Gamma(d/2)}\quad(d=D-1),
\end{equation*}
($\Omega_{d-1}$ the area of the unit sphere in $d$ dimensions). For $D=4$: $\nabla^2\Phi_{\rm N}=4\pi G\rho$ and $\Phi_{\rm N}=-GM/r$. For $D=8$, $d=7$: $\Omega_6=16\pi^3/15$, $(D-3)/(D-2)=5/6$, and a point mass $M$ has
\begin{equation*}
\Phi_8(r)=-\frac56\,\frac{\kappa M}{5\,\Omega_6\,r^5}=-\frac{5\kappa M}{32\pi^3r^5}\qquad(\kappa=8\pi G_8).
\end{equation*}
\textbf{Step 2: images.} On the torus the mass at $\vec y=0$ also sits, as far as the covering space $\mathbb R^4$ knows, at every lattice point $\ell_\perp\vec n$ (we set $b_0=1$, $L_c=\ell_\perp$). A second mass on the brane, at the distance $r$ along the brane, feels all the images:
\begin{equation*}
\Phi_{\rm N}(r)=-\frac{5\kappa M}{32\pi^3}\sum_{\vec n\in\mathbb Z^4}\frac1{(r^2+\ell_\perp^2|\vec n|^2)^{5/2}}=-\frac{5\kappa M}{32\pi^3\ell_\perp^5}\,S\Big(\frac r{\ell_\perp}\Big),\qquad S(r)=\sum_{\vec n\in\mathbb Z^4}\frac1{(r^2+|\vec n|^2)^{5/2}} .
\end{equation*}
\textbf{Step 3: Poisson summation.} For a smooth function $f$ on $\mathbb R^4$ that decays fast enough, the sum over the lattice equals the sum of its Fourier transform over the dual lattice:
\begin{equation*}
\sum_{\vec n\in\mathbb Z^4}f(\vec n)=\sum_{\vec m\in\mathbb Z^4}\hat f(2\pi\vec m),\qquad \hat f(\vec p)=\int d^4y\;e^{-i\vec p\cdot\vec y}f(\vec y).
\end{equation*}
The four-dimensional Fourier transform of $(r^2+|\vec y|^2)^{-5/2}$ is
\begin{equation*}
\int d^4y\;\frac{e^{-i\vec p\cdot\vec y}}{(r^2+|\vec y|^2)^{5/2}}=\frac{4\pi^2}p\int_0^\infty\frac{y^2J_1(py)}{(r^2+y^2)^{5/2}}\,dy=\frac{4\pi^2}{3r}\,e^{-pr}
\end{equation*}
The first step is the Fourier transform of a radial function in four dimensions ($J_1$ the Bessel function); the second uses the standard integral
\begin{equation*}
\int_0^\infty\frac{y^{\nu+1}J_\nu(py)}{(y^2+r^2)^{\mu+1}}\,dy=\frac{p^\mu r^{\nu-\mu}}{2^\mu\,\Gamma(\mu+1)}\,K_{\nu-\mu}(pr)\qquad(\nu=1,\ \mu=\tfrac32),\qquad K_{-1/2}(z)=K_{1/2}(z)=\sqrt{\frac\pi{2z}}\,e^{-z},
\end{equation*}
($K$ the modified Bessel function), which gives $\frac{p}{3r}e^{-pr}$ for the radial integral; the Python program checks the whole transform numerically to $10^{-20}$ for three pairs $(p,r)$. Hence the IDENTITY
\begin{equation*}
\boxed{\;\sum_{\vec n\in\mathbb Z^4}\frac1{(r^2+|\vec n|^2)^{5/2}}=\frac{4\pi^2}{3r}\sum_{\vec m\in\mathbb Z^4}e^{-2\pi r|\vec m|}\;}
\end{equation*}
— a sum over IMAGES (left) equals a sum over KALUZA–KLEIN MODES (right). Section 25.4 of 7+1.nb evaluates both sides (the left with the shells $r_4(k)$ up to $k=1.5\times10^6$ and a continuum tail, the right with shells up to $k=20000$):

{\footnotesize
\begin{longtable}{>{\raggedright\arraybackslash}p{\dimexpr 0.1600\linewidth-2\tabcolsep\relax}>{\raggedright\arraybackslash}p{\dimexpr 0.3000\linewidth-2\tabcolsep\relax}>{\raggedright\arraybackslash}p{\dimexpr 0.3000\linewidth-2\tabcolsep\relax}>{\raggedright\arraybackslash}p{\dimexpr 0.2400\linewidth-2\tabcolsep\relax}}
\toprule
\textbf{$r/\ell_\perp$} & \textbf{images (left side)} & \textbf{Kaluza–Klein modes (right side)} & \textbf{relative difference} \\
\midrule\endfirsthead
\toprule
\textbf{$r/\ell_\perp$} & \textbf{images (left side)} & \textbf{Kaluza–Klein modes (right side)} & \textbf{relative difference} \\
\midrule\endhead
\endfoot
\bottomrule\endlastfoot
$1/4$ & 1046.849375970 & 1046.849375970 & $-7.3\times10^{-13}$ \\
$1/2$ & 51.12525179526 & 51.12525179603 & $-1.5\times10^{-11}$ \\
$1$ & 13.40960454141 & 13.40960454217 & $-5.7\times10^{-11}$ \\
$2$ & 6.579922931695 & 6.579922932459 & $-1.2\times10^{-10}$ \\*
\end{longtable}}
\textbf{Step 4: what it means.} Restoring $\ell_\perp$ and using $8\pi G_4=\kappa/V_\perp$, $V_\perp=\ell_\perp^4$:
\begin{equation*}
\Phi_{\rm N}(r)=-\frac{5\kappa M}{32\pi^3}\cdot\frac{4\pi^2}{3V_\perp r}\sum_{\vec m}e^{-\mu_{\vec m}r}=-\frac53\,\frac{G_4M}{r}\sum_{\vec m\in\mathbb Z^4}e^{-\mu_{\vec m}r}=-\frac53\,\frac{G_4M}{r}\Big(1+8e^{-2\pi r/\ell_\perp}+24e^{-2\sqrt2\pi r/\ell_\perp}+\cdots\Big),
\end{equation*}
with $\mu_{\vec m}=2\pi|\vec m|/\ell_\perp$ and $\frac{5}{32\pi^3}\cdot\frac{4\pi^2}3\cdot8\pi=\frac53$. Every Kaluza–Klein level adds a YUKAWA potential $e^{-\mu r}/r$ of the same strength as the zero mode; the multiplicities are $r_4(k)$ (8, 24, …).

\begin{itemize}
\item \textbf{Far from the source} ($r\gg\ell_\perp$) only the zero mode survives: $\Phi_{\rm N}\to-\frac53G_4M/r$, a four-dimensional $1/r$ law. Check: $\big(\frac{3r}{4\pi^2}S(r)-1\big)\big/\big(8e^{-2\pi r}\big)=1.272$, 1.017, 1.001 at $r/\ell_\perp=1$, 2, 3 — approaching 1 as the eight lightest levels take over.
\item \textbf{Near the source} ($r\ll\ell_\perp$) the nearest image dominates and the eight-dimensional law returns, $S(r)=r^{-5}+Z(5/2)+O(r^2)$ with the Epstein value $Z(5/2)=\sum_{\vec n\ne0}|\vec n|^{-5}=\sum_kr_4(k)k^{-5/2}=8(1-4^{-3/2})\zeta(\tfrac52)\zeta(\tfrac32)=7\zeta(\tfrac52)\zeta(\tfrac32)=24.531278$ (shell sum: 24.53127769). Check: $(r^5S(r)-1)/(r^5Z(5/2))=0.99703$ at $r/\ell_\perp=1/20$ and 0.98823 at $1/10$: at a tenth of the torus size the force differs from the pure eight-dimensional one by $r^5Z(5/2)=0.00025$.
\end{itemize}
\Needspace{8\baselineskip}\subsection*{6.4 Graviton and radion: the factor \texorpdfstring{$\frac53$}{53}, and what the stabilizer changes}\addcontentsline{toc}{subsection}{6.4 Graviton and radion: the factor \texorpdfstring{$\frac53$}{53}, and what the stabilizer changes}
Why $\frac53G_4$ and not $G_4$? Without the stabilizing potential the size of the torus is a massless field, the RADION, and it is exchanged between the two masses together with the four-dimensional graviton. In 7+1.nb (section 20.3) the radion $\beta=\ln(b/b_0)$ has the canonical form $\phi=\sqrt{12}\,M_{\rm Pl}\beta$ and couples to brane matter through the trace of its energy–momentum tensor,
\begin{equation*}
\mathcal L_{\phi W}=\frac{\phi}{\sqrt3\,M_{\rm Pl}}\,T^\mu{}_\mu .
\end{equation*}
For two static masses $T^\mu{}_\mu=M\delta^3(\vec x)$; the static field equation $(-\nabla^2+m_r^2)\phi=M\delta^3/(\sqrt3M_{\rm Pl})$ gives $\phi=Me^{-m_rr}/(4\pi\sqrt3M_{\rm Pl}r)$, and the second mass feels the potential energy
\begin{equation*}
V_\phi(r)=-\frac{\phi\,M'}{\sqrt3M_{\rm Pl}}=-\frac{MM'}{12\pi M_{\rm Pl}^2}\,\frac{e^{-m_rr}}r=-\frac23\,\frac{G_4MM'}{r}\,e^{-m_rr},\qquad \frac1{12\pi M_{\rm Pl}^2}=\frac23\cdot\frac1{8\pi M_{\rm Pl}^2}=\frac23G_4 .
\end{equation*}
With $m_r=0$: graviton ($1$) $+$ radion ($\frac23$) $=\frac53$ — exactly the coefficient of the image sum. This checks two results of 7+1.nb at once: $8\pi G_4=\kappa/V_\perp$ and the normalization $1/(\sqrt3M_{\rm Pl})$ of the radion coupling. \textbf{With the stabilizing potential} the radion has the mass $m_r$ (10 GeV, a range $1/m_r=0.1$ GeV$^{-1}=0.02$ fm, far below $\ell_\perp$), and at distances $r\gg\ell_\perp,\,1/m_r$ only the graviton remains: Newton’s law $V=-G_4MM'/r$ with $G_4=G_8/V_\perp$. At $r\ll1/m_r$ (where the potential is negligible) the eight-dimensional law of step 1 holds; between $1/m_r$ and $\ell_\perp$ the force is that of eight-dimensional gravity with a stabilized volume, which 7+1.nb does not compute.

\Needspace{8\baselineskip}\subsection*{6.5 If the size changes: the radion}\addcontentsline{toc}{subsection}{6.5 If the size changes: the radion}
Since $G_4=\kappa/(8\pi b^4L_c^4)$, a changing torus changes the strength of gravity on the brane: $\dot G_4/G_4=-4H_b$ with $H_b=b'/b$. 7+1.nb holds the size with the STABILIZING POTENTIAL (an energy density in the bulk that depends on the size)
\begin{equation*}
U(b)=U_0+\frac{K_U}2\,\ln^2\frac b{b_0},\qquad \beta=\ln\frac b{b_0},\qquad m_r^2=\frac{\kappa K_U}{12},\qquad \phi=\sqrt{12}\,M_{\rm Pl}\,\beta,
\end{equation*}
whose minimum $b_0$ is the size today; $m_r$ is the radion mass. In Einstein gravity the exact equation of the radion is
\begin{equation*}
\beta''+3H_a\beta'+4\beta'^2=\kappa\Big(\frac\varepsilon6-\frac{p_a}2+\frac{p_b}3\Big)+\frac\Lambda3,
\end{equation*}
with $\varepsilon$ the energy density and $p_a$, $p_b$ the pressures along ordinary and auxiliary space of everything in the bulk (the potential contributes $\varepsilon=U$, $p_a=-U$, $p_b=-U-\frac b4U'(b)$, which gives the restoring term $-m_r^2\beta$), and $\Lambda$ the bulk cosmological constant. Brane fields see the metric $h=(b_0/b)^4g_E$, with $g_E$ the four-dimensional Einstein-frame metric; this is the origin of the coupling $\phi\,T^\mu{}_\mu/(\sqrt3M_{\rm Pl})$ of section 6.4.

\Needspace{12\baselineskip}\section*{7. The life of the torus: its size in the three regimes of the model}\addcontentsline{toc}{section}{7. The life of the torus: its size in the three regimes of the model}
7+1.nb solves the coupled field equations of the model (Einstein–Lovelock gravity with the couplings $\alpha_2$ (Gauss–Bonnet) and $\alpha_3$ (cubic), dirac16complex, the stabilizing potential, the brane) in three regimes (sections 14, 21–23). Here is what the torus does in each; $\tau$ is the time since the singularity, $H_a=a'/a$, $H_b=b'/b$.

\Needspace{8\baselineskip}\subsection*{7.1 The primordial regime: the torus contracts}\addcontentsline{toc}{subsection}{7.1 The primordial regime: the torus contracts}
Near $x_0=0$ power laws $a\propto\tau^p$, $b\propto\tau^q$ solve the vacuum equations of each Lovelock order. In Einstein gravity $3p+4q=1$ and $3p^2+4q^2=1$, so
\begin{equation*}
p_1=\frac{1+2\sqrt2}7=0.546918,\qquad q_1=\frac{2-3\sqrt2}{14}=-0.160189 .
\end{equation*}
So \textbf{ordinary space expands while the torus contracts} — a dynamical compactification. In the Gauss–Bonnet regime ($3p+4q=3$) $p_2=1.56115$, $q_2=-0.420861$; closest to the singularity, in the cubic regime, ordinary space expands like $a\propto\tau$ while the torus contracts at the finite rate $H_b\to h_0=-\sqrt{\alpha_2/(30\alpha_3)}=-182.574$ (for $\alpha_3/\alpha_2^2=10^{-6}$, in units $\alpha_2=1$); with a gas of dirac16complex, $H_b\to h_*=-\sqrt{\alpha_2/(54\alpha_3)}=-136.083$ (section 21.3 finds $-136.083$ numerically).

\Needspace{8\baselineskip}\subsection*{7.2 Without a stabilizer the torus re-expands}\addcontentsline{toc}{subsection}{7.2 Without a stabilizer the torus re-expands}
Once a gas of dirac16complex dominates, Einstein gravity with a fluid of pressures $p_a=w_a\varepsilon$, $p_b=w_b\varepsilon$ has power laws $a\propto\tau^P$, $b\propto\tau^Q$ with $3P(1+w_a)+4Q(1+w_b)=2$. The momenta along the torus are blue-shifted by its contraction, so the gas ends with $(w_a,w_b)=(0,\frac14)$ and
\begin{equation*}
(P,Q)=\Big(0,\frac25\Big):\qquad b\propto\tau^{2/5},
\end{equation*}
ordinary space stops expanding and the torus GROWS (section 21.3 finds $\tau H_b=0.399995$ and $p_b/\varepsilon=0.249999$ at $H_a=10^{-14}$). Without the stabilizing potential the 7+1 universe would decompactify.

\Needspace{8\baselineskip}\subsection*{7.3 The stabilizer catches the torus}\addcontentsline{toc}{subsection}{7.3 The stabilizer catches the torus}
With $U(b)$ the radion is caught when the expansion rate falls below $m_r$ and then oscillates about $\beta=0$ like a four-dimensional massive field in a matter-dominated universe:
\begin{equation*}
\beta\simeq\frac{A\,\sin(m_r\tau+\vartheta)}{m_r(\tau-\tau_E)},\qquad H_E=H_a+2H_b\simeq\frac2{3(\tau-\tau_E)} .
\end{equation*}
Section 22.1 integrates the coupled equations (units $\ell_2=\sqrt{\alpha_2}=1$, $m_r=10^{-3}$, start at $\tau=56.750$ with $\beta=0.4596$) and fits $A=0.4693$, $\vartheta=2.136$, $\tau_E=-1564.36$ with a largest residual of 0.013 of the amplitude. After the capture the torus is static on average, and the gas of Kaluza–Klein modes (its frozen momenta now masses) is four-dimensional matter.

\Needspace{8\baselineskip}\subsection*{7.4 Today: frozen}\addcontentsline{toc}{subsection}{7.4 Today: frozen}
With the torus static the brane obeys the four-dimensional Friedmann equations with $G_4=G_8/V_\perp$. The radion sits at a displaced minimum,
\begin{equation*}
\beta_{\rm eq}=\frac{\Lambda_4+\kappa f_W/V_\perp+\frac\kappa2(\varepsilon-3p_a)+\kappa p_b-12\alpha_2H_a^2(H_a'+H_a^2)}{3m_r^2},
\end{equation*}
($\Lambda_4$ the four-dimensional cosmological constant, $f_W$ the brane tension) — numerically $\beta_{\rm eq}=1.7\times10^{-86}$ from $\Lambda_4$ and matter, $5.6\times10^{-24}$ from a tension $f_W=M_8^4$ (section 22.4, $m_r=10$ GeV): $G_4$ is constant to that accuracy, as experiments require. The Gauss–Bonnet and cubic corrections are $\alpha_2H_0^2=2.1\times10^{-92}$ and $\alpha_3H_0^4=4.3\times10^{-190}$.

\Needspace{8\baselineskip}\subsection*{7.5 The far future: frozen forever (or inflating without the stabilizer)}\addcontentsline{toc}{subsection}{7.5 The far future: frozen forever (or inflating without the stabilizer)}
In the de Sitter future ($H_\infty=H_0\Omega_\Lambda^{1/2}=1.808\times10^{-18}$ s$^{-1}$) the radion relaxes to $3m_r^2\beta_{\rm eq}=3H^2-12\alpha_2H^4$, of order $H_\infty^2/m_r^2$ (section 23.2; a numerical test in units $\ell_2=1$ reaches $\beta_{\rm eq}=0.0102041$). Without the stabilizer the far future would instead be eight-dimensional de Sitter space,
\begin{equation*}
a=b=e^{Hx_0},\qquad 21H^2-420\alpha_2H^4+2520\alpha_3H^6=\Lambda,
\end{equation*}
in which the torus inflates and $G_4\propto b^{-4}$ falls like $e^{-4Hx_0}$ — excluded by the observed constancy of $G$.

\Needspace{12\baselineskip}\section*{8. The torus and the world brane: Kaluza–Klein momentum}\addcontentsline{toc}{section}{8. The torus and the world brane: Kaluza–Klein momentum}
\Needspace{8\baselineskip}\subsection*{8.1 Conserved in the bulk}\addcontentsline{toc}{subsection}{8.1 Conserved in the bulk}
The translations of the torus are symmetries of the bulk, so their charge, the Kaluza–Klein momentum $\vec n$, is conserved by every local interaction in the bulk. \textbf{Example 8.1:} the overlap of three modes, which is the amplitude of a cubic interaction,
\begin{equation*}
\int_Fb^4d^4y\;\phi_{\vec n}\,\phi_{\vec m}\,\phi^*_{\vec l}=\frac{\delta_{\vec n+\vec m,\,\vec l}}{\sqrt{V_\perp}}
\end{equation*}
is $1/\sqrt{V_\perp}$ for $\vec n=(1,0,0,2)$, $\vec m=(0,1,-1,0)$, $\vec l=(1,1,-1,2)=\vec n+\vec m$ and 0 for $\vec l=(1,1,-1,1)$ (section 25.5).

\Needspace{8\baselineskip}\subsection*{8.2 Not conserved at the brane}\addcontentsline{toc}{subsection}{8.2 Not conserved at the brane}
The world brane sits at the single point $\vec y=0$ and is not invariant under the translations. At the brane every mode has the same value, $\phi_{\vec n}(0)=1/\sqrt{V_\perp}$, so a bulk field seen from the brane is
\begin{equation*}
\psi(x,\vec y=0)=\frac1{\sqrt{V_\perp}}\sum_{\vec n\in\mathbb Z^4}\chi_{\vec n}(x),\qquad \big(\phi_{\vec n}\phi_{\vec m}\phi^*_{\vec l}\big)(0)=V_\perp^{-3/2}\ \text{ for all }\vec n,\vec m,\vec l,
\end{equation*}
($\chi_{\vec n}$ the four-dimensional field of the level $\vec n$): an interaction on the brane couples EVERY level with the same strength and can change $\vec n$ by any amount. The brane couplings of 7+1.nb (the portals) therefore act on the whole tower at once — the reason why the Higgs portal of section 24 couples every pair of levels.

\Needspace{8\baselineskip}\subsection*{8.3 The delta function of the torus, and smeared sources}\addcontentsline{toc}{subsection}{8.3 The delta function of the torus, and smeared sources}
The delta function of the torus contains every mode with the same weight,
\begin{equation*}
\delta^4(\vec y)=\frac1{L_c^4}\sum_{\vec n\in\mathbb Z^4}e^{i\vec q_{\vec n}\cdot\vec y},\qquad \int_{-L_c/2}^{L_c/2}e^{-iq_ny}\,\delta(y)\,dy=1,
\end{equation*}
and its average over the torus (its zero mode) is $1/L_c^4$. A source on the brane with the four-dimensional energy density $\rho_4$ has the eight-dimensional density $\rho_4\,\delta^4(\vec y)/b^4$, whose average over the torus is
\begin{equation*}
\varepsilon_8=\frac1{V_\perp}\int_Fb^4d^4y\;\frac{\rho_4\,\delta^4(\vec y)}{b^4}=\frac{\rho_4}{V_\perp} :
\end{equation*}
this SMEARING is how the homogeneous equations of 7+1.nb treat the brane (a brane tension $f_W$ gives $\varepsilon=f_W/V_\perp$, $p_a=-f_W/V_\perp$, $p_b=0$).

\Needspace{12\baselineskip}\section*{9. The vacuum energy of the torus}\addcontentsline{toc}{section}{9. The vacuum energy of the torus}
Quantum fields on the torus have a zero-point energy that depends on its size — a CASIMIR energy. For a field of mass $m$ with periodic conditions, per real bosonic degree of freedom and per unit volume of the brane (section 24.5 of 7+1.nb),
\begin{align*}
\rho_{\rm Cas}^{(1)}&=-\frac2{\ell_\perp^4}\sum_{\vec n\ne0}\Big(\frac{m\ell_\perp}{2\pi|\vec n|}\Big)^4K_4(m\ell_\perp|\vec n|)\;\xrightarrow{\;m\to0\;}\;-\frac{\Gamma(4)}{\pi^4\ell_\perp^4}\,Z(4), \\
Z(4)&=\sum_{\vec n\ne0}\frac1{|\vec n|^8}=\sum_{k\ge1}\frac{r_4(k)}{k^4}=8(1-4^{-3})\,\zeta(4)\,\zeta(3),
\end{align*}
($K_4$ the modified Bessel function, $\zeta$ the Riemann zeta function); a fermion contributes with the opposite sign, and dirac16complex has 16 real degrees of freedom on shell. With $\ell_\perp$ of section 6.2 and the observed dark-energy density $\rho_\Lambda=2.518\times10^{-47}$ GeV$^4$, a massless dirac16complex gives $6.76\times10^{34}\,\rho_\Lambda$ and the graviton (20 degrees of freedom) $-8.45\times10^{34}\,\rho_\Lambda$; the dirac16complex term equals $\rho_\Lambda$ only for $m\ell_\perp=92.23$ ($m=59.10$ MeV). While the torus is frozen these energies are constant ($w=-1$): they are terms of the four-dimensional cosmological constant, not a dynamical dark energy.

\Needspace{12\baselineskip}\section*{10. The torus compared with a 4-sphere}\addcontentsline{toc}{section}{10. The torus compared with a 4-sphere}
Why is the auxiliary space of 7+1 a torus? The clearest way to see what the choice implies is to replace it by the other simple compact four-dimensional space, the 4-SPHERE $S^4$ of radius $r$, and compare (section 26 of 7+1.nb). On $S^4$ the coordinates $x_4,x_5,x_6\in[0,\pi]$, $x_7\in[0,2\pi)$ are angles, and
\begin{equation*}
g_{kl}=-r^2\hat\gamma_{kl},\qquad \hat\gamma=\mathrm{diag}\big(1,\sin^2x_4,\sin^2x_4\sin^2x_5,\sin^2x_4\sin^2x_5\sin^2x_6\big),\qquad V=\frac{8\pi^2}3r^4,\qquad R_{kl}=3\hat\gamma_{kl} .
\end{equation*}
\Needspace{8\baselineskip}\subsection*{10.1 Flat versus curved: the size as a field}\addcontentsline{toc}{subsection}{10.1 Flat versus curved: the size as a field}
Let the size depend on the point of ordinary spacetime, $b=b_0e^{\beta(x)}$ (torus) or $r=r_0e^{\beta(x)}$ (sphere), with the four-dimensional metric rescaled so that Newton’s constant stays constant (the Einstein frame):
\begin{equation*}
ds^2=e^{-4\beta(x)}g_{E\mu\nu}dx^\mu dx^\nu-e^{2\beta(x)}\gamma_{kl}dy^kdy^l,\qquad \gamma=b_0^2\delta\ \ (T^4),\quad \gamma=r_0^2\hat\gamma\ \ (S^4).
\end{equation*}
Integrated over the extra space, the eight-dimensional action becomes $\sqrt{|g_E|}[-\frac{M_{\rm Pl}^2}2R_E+6M_{\rm Pl}^2(\partial\beta)^2-V(\beta)]$ with $M_{\rm Pl}^2=V_0/\kappa$, so $\phi=\sqrt{12}M_{\rm Pl}\beta$ is the canonical radion on both spaces, and
\begin{equation*}
V(\beta)=M_{\rm Pl}^2\Big[e^{-4\beta}\big(\Lambda+\kappa U(\beta)\big)-\frac{6\sigma}{r_0^2}e^{-6\beta}\Big],\qquad \sigma=0\ (T^4),\quad \sigma=1\ (S^4)
\end{equation*}
($\Lambda$ the bulk cosmological constant, $U$ the stabilizing energy density). The flat torus is a MODULUS: for $\Lambda=U=0$ every size is a solution. The curved sphere is not: its curvature energy $-6M_{\rm Pl}^2e^{-6\beta}/r_0^2$ drives it to collapse. With $U=U_0+U_1\beta+\frac{K_U}2\beta^2$ a flat four-dimensional vacuum ($V(0)=V'(0)=0$) needs
\begin{equation*}
T^4:\ \Lambda+\kappa U_0=0,\ U_1=0,\ m_r^2=\frac{\kappa K_U}{12};\qquad S^4:\ (\Lambda+\kappa U_0)r_0^2=6,\ \kappa U_1r_0^2=-12,\ m_r^2=\frac{\kappa K_U}{12}-2\,\frac1{r_0^2},
\end{equation*}
($m_r^2=V''(0)/(12M_{\rm Pl}^2)$). A sphere held by $\Lambda$ alone ($dS_4\times S^4$, $\Lambda r_0^2=9$) is unstable, $m^2=-6\,H^2$.

\Needspace{8\baselineskip}\subsection*{10.2 Levels and their counting}\addcontentsline{toc}{subsection}{10.2 Levels and their counting}
{\footnotesize
\begin{longtable}{>{\raggedright\arraybackslash}p{\dimexpr 0.2400\linewidth-2\tabcolsep\relax}>{\raggedright\arraybackslash}p{\dimexpr 0.3800\linewidth-2\tabcolsep\relax}>{\raggedright\arraybackslash}p{\dimexpr 0.3800\linewidth-2\tabcolsep\relax}}
\toprule
\textbf{} & \textbf{$T^4$} & \textbf{$S^4$} \\
\midrule\endfirsthead
\toprule
\textbf{} & \textbf{$T^4$} & \textbf{$S^4$} \\
\midrule\endhead
\endfoot
\bottomrule\endlastfoot
eigenfunctions of $-\Delta$ & $e^{2\pi i\vec n\cdot\vec y/\ell_\perp}/\sqrt V$ & spherical harmonics of degree $l$ \\
masses & $2\pi|\vec n|/\ell_\perp$ & $\sqrt{l(l+3)}/r$ \\
multiplicity & $r_4(k)=8\sigma(k)-32\sigma(k/4)$ & $d_l=\frac{(2l+3)(l+1)(l+2)}6$ \\
first multiplicities & $1,8,24,32,24,\dots$ ($k=0,1,2,3,4$) & $1,5,14,30,55,91,\dots$ ($l=0,\dots,5$) \\
Weyl’s law & $N(\mu)\simeq V\mu^4/(32\pi^2)$ & the same; exactly $N=y^2(y^2-1)/12$, $y=L+2$ \\*
\end{longtable}}
Both obey Weyl’s law, which depends only on the volume: the number of heavy levels cannot tell the two spaces apart. \textbf{Example 10.2:} with the same volume $V_\perp$ (the same $G_4$), the sphere has the radius $r_{S^4}=(3V_\perp/(8\pi^2))^{1/4}=1.3595\times10^{-13}$ m; below 5 GeV the torus has 11733511217161 states (lattice count), the sphere 11737500461050 (exact), Weyl’s law gives $1.17335\times10^{13}$. The LIGHTEST levels differ: the torus at $2\pi/\ell_\perp=4.0265$ MeV (8 states), the sphere at $2/r_{S^4}=2.9030$ MeV (5 states).

\Needspace{8\baselineskip}\subsection*{10.3 What a brane sees}\addcontentsline{toc}{subsection}{10.3 What a brane sees}
A brane at $\vec y=0$ of the torus sees every mode with the value $1/\sqrt V$: per level it couples to one combination of the $r_4(k)$ modes with the strength $\sqrt{r_4(k)/V}$. A brane at the north pole $x_4=0$ of the sphere sees only the ZONAL harmonic of each level,
\begin{equation*}
Z_l=\sqrt{\frac{d_l}V}\,\frac{C_l^{(3/2)}(\cos x_4)}{C_l^{(3/2)}(1)},\qquad Z_l(0)=\sqrt{\frac{d_l}V}
\end{equation*}
($C_l^{(3/2)}$ the Gegenbauer polynomial): one state per level with the strength $\sqrt{d_l/V}$. The brane keeps 384 symmetries of the torus and 6 of the 10 rotations of the sphere ($SO(4)\subset SO(5)$). On both spaces the brane sees the same induced metric $e^{-4\beta}g_E$, so the radion couples to brane matter in the same way, $\phi T^\mu{}_\mu/(\sqrt3M_{\rm Pl})$.

\Needspace{8\baselineskip}\subsection*{10.4 Spinors: zero modes only on the torus}\addcontentsline{toc}{subsection}{10.4 Spinors: zero modes only on the torus}
With periodic spinors the torus has the constant spinors as zero modes of its Dirac operator: dirac16complex gives four massless brane Dirac fermions for $m=0$ (section 4.5). The sphere has none: the eigenvalues of its Dirac operator are $\pm i(l+2)/r$ with the degeneracy $4\binom{l+3}3$ for each sign, the lowest belonging to the Killing spinors, which reach Friedrich’s bound $|\lambda|^2\ge nR/(4(n-1))=4/r^2$ ($R=12/r^2$). On $S^4$ the lightest dirac16complex states are eight Dirac fermions of mass $\sqrt{m^2+4/r^2}$, at least 2.9030 MeV for $r=r_{S^4}$.

\Needspace{8\baselineskip}\subsection*{10.5 Side by side}\addcontentsline{toc}{subsection}{10.5 Side by side}
{\footnotesize
\begin{longtable}{>{\raggedright\arraybackslash}p{\dimexpr 0.3000\linewidth-2\tabcolsep\relax}>{\raggedright\arraybackslash}p{\dimexpr 0.3500\linewidth-2\tabcolsep\relax}>{\raggedright\arraybackslash}p{\dimexpr 0.3500\linewidth-2\tabcolsep\relax}}
\toprule
\textbf{property} & \textbf{$T^4$ (7+1.nb)} & \textbf{$S^4$ (comparison)} \\
\midrule\endfirsthead
\toprule
\textbf{property} & \textbf{$T^4$ (7+1.nb)} & \textbf{$S^4$ (comparison)} \\
\midrule\endhead
\endfoot
\bottomrule\endlastfoot
curvature & flat & $R_{kl}=3\hat\gamma_{kl}$, $R=12/r^2$ \\
volume & $b^4L_c^4$ & $\frac{8\pi^2}3r^4$ \\
size without stabilizer & modulus (any size) & collapses \\
stabilization & $\Lambda+\kappa U_0=0$, $U_1=0$ & $(\Lambda+\kappa U_0)r_0^2=6$, $\kappa U_1r_0^2=-12$ \\
radion mass & $\kappa K_U/12$ & $\kappa K_U/12-2/r_0^2$ \\
size today ($V=V_\perp$) & $\ell_\perp=3.079\times10^{-13}$ m & $r_{S^4}=1.359\times10^{-13}$ m \\
lightest level & 4.027 MeV, 8 states & 2.903 MeV, 5 states \\
symmetry kept by a brane & 384 signed permutations & $SO(4)$ \\
Dirac zero modes & 4 & none \\*
\end{longtable}}
\Needspace{12\baselineskip}\section*{11. Validity and limitations}\addcontentsline{toc}{section}{11. Validity and limitations}
\begin{itemize}
\item The torus of 7+1.nb is the CUBIC torus: four circles of the same period, at right angles. A general flat torus $\mathbb R^4/\Lambda$ with any lattice $\Lambda$ would have the masses $\mu^2=\vec q\cdot G^{-1}\vec q$ ($G$ the metric of the lattice, $\vec q$ in the dual lattice); its shape is not dynamical in 7+1.nb.
\item Only the volume is stabilized ($U(b)$); the nine shape moduli and the four graviphotons (section 4.7) are not discussed in 7+1.nb.
\item Fermions are periodic around all four circles (section 4.6).
\item The 4-sphere of section 10 is a comparison worked out in section 26 of 7+1.nb; the model itself uses the torus.
\item The Newton-law example (section 6.3) is linearized gravity of a point mass without the stabilizing potential; the stabilizer is added through the radion mass only (section 6.4). The gravitational field of the brane itself is not included.
\item The numbers today assume $M_8=10$ TeV (and $M_{\rm Pl}=2.435\times10^{18}$ GeV); $\ell_\perp\propto M_8^{-3/2}$, so a different $M_8$ changes them accordingly.
\item Continuum formulas such as $N(\mu)=V_\perp\mu^4/(32\pi^2)$ hold for $\mu\gg2\pi/\ell_\perp$; their accuracy is that of Weyl’s law (section 5.2).
\end{itemize}
\Needspace{12\baselineskip}\section*{12. Verification}\addcontentsline{toc}{section}{12. Verification}
Every worked example of this document is a check of section 25 of 7+1.nb (Wolfram Language 15.0.1; evaluated in a fresh kernel, 0 message cells) and of the independent Python program tools/ref\_torus\_brane.py (sympy, numpy, mpmath; its own lattice sums, quadratures and Christoffel symbols). The checks of section 25:

{\footnotesize
\begin{longtable}{>{\raggedright\arraybackslash}p{\dimexpr 0.0500\linewidth-2\tabcolsep\relax}>{\raggedright\arraybackslash}p{\dimexpr 0.7700\linewidth-2\tabcolsep\relax}>{\raggedright\arraybackslash}p{\dimexpr 0.0800\linewidth-2\tabcolsep\relax}>{\raggedright\arraybackslash}p{\dimexpr 0.1000\linewidth-2\tabcolsep\relax}}
\toprule
\textbf{\#} & \textbf{check} & \textbf{result} & \textbf{section} \\
\midrule\endfirsthead
\toprule
\textbf{\#} & \textbf{check} & \textbf{result} & \textbf{section} \\
\midrule\endhead
\endfoot
\bottomrule\endlastfoot
1 & \texttt{auxiliary torus: for the cosmological metric det(-g\_\allowbreak{}kl) == b\textasciicircum{}8 (k, l = 4..7), and the proper volume of the fundamental domain [0, Lc)\textasciicircum{}4, Integrate[Sqrt[det(-g\_\allowbreak{}kl)], \{y, 0, Lc\}\textasciicircum{}4], is b\textasciicircum{}4 Lc\textasciicircum{}4 = V\_\allowbreak{}perp} & True & 25.1 \\
2 & \texttt{only the product b Lc of the auxiliary torus is physical: y = lambda y' maps the torus of period Lc and size b to the one of period Lc/\allowbreak{}lambda and size lambda b; the line element b\textasciicircum{}2 dy.dy, the proper period b Lc and V\_\allowbreak{}perp = b\textasciicircum{}4 Lc\textasciicircum{}4 are unchanged} & True & 25.1 \\
3 & \texttt{symmetries of the auxiliary torus that fix y = 0 (and the brane): the 384 = 2\textasciicircum{}4 4! signed permutation matrices are distinct orthogonal integer matrices (they map Z\textasciicircum{}4 onto itself and preserve |n|\textasciicircum{}2); the shells |n|\textasciicircum{}2 = 1, 2, 3, 4, 9 split into orbits of sizes \{8\}, \{24\}, \{32\}, \{8, 16\}, \{8, 96\}} & True & 25.1 \\
4 & \texttt{Fourier modes of the auxiliary torus: phi\_\allowbreak{}n = Exp[I q\_\allowbreak{}n.y]/\allowbreak{}Sqrt[V\_\allowbreak{}perp] with q\_\allowbreak{}n = 2 Pi n/\allowbreak{}Lc (n in Z\textasciicircum{}4) are orthonormal for the proper volume element b\textasciicircum{}4 d\textasciicircum{}4y (exact integrals: n = m = (1, 0, -2, 3); n = (1, 0, -2, 3), m = (1, 1, -2, 3); n = 0, m = (0, 2, 0, 0))} & True & 25.2 \\
5 & \texttt{Parseval on the auxiliary torus: f(y) = Prod\_\allowbreak{}k cos\textasciicircum{}2(Pi y\_\allowbreak{}k/\allowbreak{}Lc) = Prod\_\allowbreak{}k (1/\allowbreak{}2 + (1/\allowbreak{}4) e\textasciicircum{}(2 Pi I y\_\allowbreak{}k/\allowbreak{}Lc) + (1/\allowbreak{}4) e\textasciicircum{}(-2 Pi I y\_\allowbreak{}k/\allowbreak{}Lc)) has the 81 non-zero coefficients c\_\allowbreak{}n = Sqrt[V\_\allowbreak{}perp] Prod\_\allowbreak{}k w(n\_\allowbreak{}k), w(0) = 1/\allowbreak{}2, w(+-1) = 1/\allowbreak{}4 (three of them checked by integration), and Sum\_\allowbreak{}n |c\_\allowbreak{}n|\textasciicircum{}2 == Int b\textasciicircum{}4 d\textasciicircum{}4y f\textasciicircum{}2 == (3/\allowbreak{}8)\textasciicircum{}4 V\_\allowbreak{}perp} & True & 25.2 \\
6 & \texttt{the Laplacian of the auxiliary torus: -(1/\allowbreak{}b\textasciicircum{}2) Sum\_\allowbreak{}k d\textasciicircum{}2 phi\_\allowbreak{}n/\allowbreak{}dy\_\allowbreak{}k\textasciicircum{}2 == (|q\_\allowbreak{}n|\textasciicircum{}2/\allowbreak{}b\textasciicircum{}2) phi\_\allowbreak{}n (n = (2, -1, 0, 3))} & True & 25.2 \\
7 & \texttt{scalar Kaluza-Klein masses: for the static metric diag(1, -1, -1, -1, -b0\textasciicircum{}2, ..., -b0\textasciicircum{}2) the Klein-Gordon operator Box + m\textasciicircum{}2 (Box f = |g|\textasciicircum{}(-1/\allowbreak{}2) d\_\allowbreak{}A(|g|\textasciicircum{}(1/\allowbreak{}2) g\textasciicircum{}AB d\_\allowbreak{}B f)) on Exp[-I E x0 + I k.x + I q.y] gives (-E\textasciicircum{}2 + |k|\textasciicircum{}2 + |q|\textasciicircum{}2/\allowbreak{}b0\textasciicircum{}2 + m\textasciicircum{}2) times the wave: in ordinary spacetime a particle of mass Sqrt[m\textasciicircum{}2 + |q|\textasciicircum{}2/\allowbreak{}b0\textasciicircum{}2]} & True & 25.2 \\
8 & \texttt{Jacobi's four-square theorem: the number of n in Z\textasciicircum{}4 with |n|\textasciicircum{}2 = k is r\_\allowbreak{}4(k) = 8 sigma(k) - 32 sigma(k/\allowbreak{}4) for k = 1..100 (direct count over \{-10, ..., 10\}\textasciicircum{}4)} & True & 25.3 \\
9 & \texttt{Weyl's law: the number N(R) of n in Z\textasciicircum{}4 with |n| <= R approaches the volume Pi\textasciicircum{}2 R\textasciicircum{}4/\allowbreak{}2 of the four-ball: N(10) = 49689, N(100) = 493490641, and |N(R)/\allowbreak{}(Pi\textasciicircum{}2 R\textasciicircum{}4/\allowbreak{}2) - 1| < 10\textasciicircum{}-2, 10\textasciicircum{}-3, 3 10\textasciicircum{}-5, 5 10\textasciicircum{}-6, 10\textasciicircum{}-6 for R = 10, 30, 100, 300, 1000} & True & 25.3 \\
10 & \texttt{the auxiliary torus today (M8 = 10 TeV, M\_\allowbreak{}Pl = 2.435 10\textasciicircum{}18 GeV): V\_\allowbreak{}perp = M\_\allowbreak{}Pl\textasciicircum{}2/\allowbreak{}M8\textasciicircum{}6 = 5.929 10\textasciicircum{}12 GeV\textasciicircum{}-4, l\_\allowbreak{}perp = V\_\allowbreak{}perp\textasciicircum{}(1/\allowbreak{}4) = b0 Lc = 1560.4 GeV\textasciicircum{}-1 = 3.079 10\textasciicircum{}-13 m, lowest Kaluza-Klein mass 2 Pi/\allowbreak{}l\_\allowbreak{}perp = 4.027 MeV; N(mu) = V\_\allowbreak{}perp mu\textasciicircum{}4/\allowbreak{}(32 Pi\textasciicircum{}2) = 1.17 10\textasciicircum{}13 levels below m\_\allowbreak{}r/\allowbreak{}2 = 5 GeV (the lattice count of section 22.2 agrees to 10\textasciicircum{}-5) and N(M8) = M\_\allowbreak{}Pl\textasciicircum{}2/\allowbreak{}(32 Pi\textasciicircum{}2 M8\textasciicircum{}2)} & True & 25.3 \\
11 & \texttt{Newton's law in eight dimensions: Omega\_\allowbreak{}6 = 2 Pi\textasciicircum{}(7/\allowbreak{}2)/\allowbreak{}Gamma(7/\allowbreak{}2) = 16 Pi\textasciicircum{}3/\allowbreak{}15, r\textasciicircum{}-5 is harmonic away from r = 0 in seven space dimensions ((1/\allowbreak{}r\textasciicircum{}6) d/\allowbreak{}dr (r\textasciicircum{}6 d/\allowbreak{}dr r\textasciicircum{}-5) = 0), and grad\textasciicircum{}2 Phi = 8 Pi G8 ((D - 3)/\allowbreak{}(D - 2)) M delta\textasciicircum{}7 (D = 8) gives Phi = -(5/\allowbreak{}6) kappa M/\allowbreak{}(5 Omega\_\allowbreak{}6 r\textasciicircum{}5) = -5 kappa M/\allowbreak{}(32 Pi\textasciicircum{}3 r\textasciicircum{}5); the same formula in D = 4 gives -G M/\allowbreak{}r} & True & 25.4 \\
12 & \texttt{images = Kaluza-Klein modes (Poisson summation; lengths in units of l\_\allowbreak{}perp): Sum\_\allowbreak{}n (r\textasciicircum{}2 + |n|\textasciicircum{}2)\textasciicircum{}(-5/\allowbreak{}2) == (4 Pi\textasciicircum{}2/\allowbreak{}(3 r)) Sum\_\allowbreak{}m Exp[-2 Pi r |m|] to 10\textasciicircum{}-8 at r = 1/\allowbreak{}4, 1/\allowbreak{}2, 1, 2 (left: shells with r\_\allowbreak{}4(k) to k = 1.5 10\textasciicircum{}6 and the continuum tail Pi\textasciicircum{}2 Int k (r\textasciicircum{}2 + k)\textasciicircum{}(-5/\allowbreak{}2) dk; right: shells to k = 20000)} & True & 25.4 \\
13 & \texttt{near and far field: Z(5/\allowbreak{}2) = Sum\_\allowbreak{}(n != 0) |n|\textasciicircum{}-5 = Sum\_\allowbreak{}k r\_\allowbreak{}4(k) k\textasciicircum{}(-5/\allowbreak{}2) = 7 zeta(5/\allowbreak{}2) zeta(3/\allowbreak{}2) = 24.531 (shell sum with tail, to 10\textasciicircum{}-8); near the source r\textasciicircum{}5 S(r) - 1 -> r\textasciicircum{}5 Z(5/\allowbreak{}2) (the ratio lies between 0.98 and 1 at r = 1/\allowbreak{}10 and 1/\allowbreak{}20, closer to 1 at 1/\allowbreak{}20): the eight-dimensional law; far from it 3 r S/\allowbreak{}(4 Pi\textasciicircum{}2) - 1 -> 8 Exp[-2 Pi r], the 8 lightest Kaluza-Klein modes (ratio 1 to 2 10\textasciicircum{}-2 at r = 2 and 3)} & True & 25.4 \\
14 & \texttt{graviton + radion: far from the source the images add up to (5 kappa M/\allowbreak{}(32 Pi\textasciicircum{}3)) (4 Pi\textasciicircum{}2/\allowbreak{}(3 V\_\allowbreak{}perp r)) == (5/\allowbreak{}3) G4 M/\allowbreak{}r with G4 = kappa/\allowbreak{}(8 Pi V\_\allowbreak{}perp); the radion of section 20.3, coupled to matter by (phi/\allowbreak{}(Sqrt[3] M\_\allowbreak{}Pl)) T, gives the potential -(1/\allowbreak{}(Sqrt[3] M\_\allowbreak{}Pl))\textasciicircum{}2 M M' Exp[-m\_\allowbreak{}r r]/\allowbreak{}(4 Pi r), and (1/\allowbreak{}(Sqrt[3] M\_\allowbreak{}Pl))\textasciicircum{}2/\allowbreak{}(4 Pi) == (2/\allowbreak{}3) G4 (8 Pi G4 = 1/\allowbreak{}M\_\allowbreak{}Pl\textasciicircum{}2): graviton 1 + massless radion 2/\allowbreak{}3 = 5/\allowbreak{}3} & True & 25.4 \\
15 & \texttt{Kaluza-Klein momentum is conserved in the bulk and not at the brane: Int b\textasciicircum{}4 d\textasciicircum{}4y phi\_\allowbreak{}n phi\_\allowbreak{}m phi\_\allowbreak{}l\textasciicircum{}* == 1/\allowbreak{}Sqrt[V\_\allowbreak{}perp] for l = n + m and 0 otherwise (exact: n = (1, 0, 0, 2), m = (0, 1, -1, 0), l = (1, 1, -1, 2) and (1, 1, -1, 1)); at y = 0, phi\_\allowbreak{}n phi\_\allowbreak{}m phi\_\allowbreak{}l\textasciicircum{}* == V\_\allowbreak{}perp\textasciicircum{}(-3/\allowbreak{}2) for every n, m, l} & True & 25.5 \\
16 & \texttt{the delta function of the torus: Int\_\allowbreak{}(-Lc/\allowbreak{}2)\textasciicircum{}(Lc/\allowbreak{}2) Exp[-I q\_\allowbreak{}n y] delta(y) dy == 1 for every integer n, so delta\textasciicircum{}4(y) = Lc\textasciicircum{}-4 Sum\_\allowbreak{}n Exp[I q\_\allowbreak{}n.y] (every mode with the weight 1/\allowbreak{}Lc\textasciicircum{}4); a brane source rho4 delta\textasciicircum{}4(y)/\allowbreak{}b\textasciicircum{}4 has Int b\textasciicircum{}4 d\textasciicircum{}4y (rho4 delta\textasciicircum{}4(y)/\allowbreak{}b\textasciicircum{}4) == rho4, i.e. the average rho4/\allowbreak{}V\_\allowbreak{}perp over the torus} & True & 25.5 \\
17 & \texttt{the world brane W = \{x4 = x5 = x6 = x7 = 0\}: the induced metric h = J\textasciicircum{}T.g.J of the embedding X = (x0, x1, x2, x3, 0, 0, 0, 0) is, for the cosmological metric of section 11.2, the Friedmann-Robertson-Walker metric diag(1, -a\textasciicircum{}2, -a\textasciicircum{}2, -a\textasciicircum{}2) (independent of b); for the inhomogeneous metric of section 11.3 it is diag(n\textasciicircum{}2, -a\textasciicircum{}2, -a\textasciicircum{}2, -a\textasciicircum{}2) at x4 = 0 with n = 1 + x0\textasciicircum{}2/\allowbreak{}11, a = 1 + x0/\allowbreak{}3} & True & 25.6 \\
18 & \texttt{brane integrals in eight-dimensional form: at the brane Sqrt|det g| == Sqrt|det h| b\textasciicircum{}4 with b\textasciicircum{}4 = Sqrt[det(-g\_\allowbreak{}kl)] (cosmological metric; inhomogeneous metric exactly at x0 = 1/\allowbreak{}2), so Int\_\allowbreak{}W d\textasciicircum{}4x Sqrt|h| L == Int d\textasciicircum{}8X Sqrt|g| (delta\textasciicircum{}4(y)/\allowbreak{}b\textasciicircum{}4) L; test: Int d\textasciicircum{}4y b\textasciicircum{}4 (delta\textasciicircum{}4(y)/\allowbreak{}b\textasciicircum{}4) F(y) == F(0) = 1 for F = Exp[x4 + 2 x5] Cos[x6 - x7]} & True & 25.6 \\
19 & \texttt{how the brane sits in the bulk: the Christoffel symbols \{k mu nu\} at the brane (k = 4..7, mu, nu = 0..3; = (1/\allowbreak{}(2 b\textasciicircum{}2)) d\_\allowbreak{}k g\_\allowbreak{}mu nu for these metrics, b times the extrinsic curvature) all vanish for the cosmological metric: the brane is totally geodesic, and a free particle moving along it stays on it. For the inhomogeneous metric at x0 = 1/\allowbreak{}2 the only non-zero one is \{4 00\} = n (d\_\allowbreak{}4 n)/\allowbreak{}b\textasciicircum{}2 = 1458/\allowbreak{}58619: a free particle at rest on the brane is accelerated off it} & True & 25.6 \\
20 & \texttt{explicit forms: every Gamma17[A] is a signed permutation matrix (one entry +-1 in each row and column), and 2 Sqrt[2] U has only the entries 0, +-1, +-I, eight non-zero ones in each row and in each column} & True & 25.7 \\
21 & \texttt{brane spinor (x) transverse spinor: Gamma17[mu].Gamma17[nu] == U.(gamma\_\allowbreak{}mu.gamma\_\allowbreak{}nu (x) 1\_\allowbreak{}4).U\textasciicircum{}dag, Gamma17[k].Gamma17[l] == U.(1\_\allowbreak{}4 (x) rho\_\allowbreak{}k.rho\_\allowbreak{}l).U\textasciicircum{}dag and Gamma17[mu].Gamma17[k] == U.(gamma\_\allowbreak{}mu.gamma5 (x) rho\_\allowbreak{}k).U\textasciicircum{}dag (mu, nu = 0..3, k, l = 4..7): the Lorentz group of the brane acts on the brane spinor index, the rotations about the brane on the transverse index, and the generators that mix them are broken by the brane} & True & 25.7 \\
22 & \texttt{the transverse index carries Spin(4) = SU(2) x SU(2): rho4567 == -gamma5 (4 x 4), it commutes with every rho\_\allowbreak{}k.rho\_\allowbreak{}l and anticommutes with every rho\_\allowbreak{}k; for each even permutation (k, l, m, n) of (4, 5, 6, 7) the combination rho\_\allowbreak{}k.rho\_\allowbreak{}l + s rho\_\allowbreak{}m.rho\_\allowbreak{}n vanishes on the eigenspace rho4567 = s, and rho\_\allowbreak{}k.rho\_\allowbreak{}l - s rho\_\allowbreak{}m.rho\_\allowbreak{}n acts there with rank 2: the four values of a form (2, 1) + (1, 2)} & True & 25.7 \\
23 & \texttt{chirality at the brane: PL17 == U.(P\_\allowbreak{}L (x) P\_\allowbreak{}+ + P\_\allowbreak{}R (x) P\_\allowbreak{}-).U\textasciicircum{}dag with P\_\allowbreak{}L,R = (1 -+ gamma5)/\allowbreak{}2 and P\_\allowbreak{}+- = (1 +- rho4567)/\allowbreak{}2: the left-handed half psi\_\allowbreak{}L (chi17 = -1, 8 components) is two brane-left-handed Weyl spinors with rho4567 = +1 and two brane-right-handed ones with rho4567 = -1 (each part of rank 4)} & True & 25.7 \\
24 & \texttt{an explicit spinor: psi = U.(u (x) e\_\allowbreak{}1) with the brane spinor u = (1, 2, 0, -1): psibar psi == ubar u == 4 and psibar Gamma\textasciicircum{}mu psi == ubar gamma\textasciicircum{}mu u (mu = 0..3)} & True & 25.7 \\
25 & \texttt{brane tension: for S\_\allowbreak{}W = -Int d\textasciicircum{}4x Sqrt|h| f\_\allowbreak{}W and h = diag(1/\allowbreak{}u0, -1/\allowbreak{}u1, -1/\allowbreak{}u2, -1/\allowbreak{}u3) (h\textasciicircum{}00 = u0, h\textasciicircum{}ii = -u\_\allowbreak{}i), T\_\allowbreak{}mu nu = (2/\allowbreak{}Sqrt|h|) dS\_\allowbreak{}W/\allowbreak{}dh\textasciicircum{}mu nu == f\_\allowbreak{}W h\_\allowbreak{}mu nu, T\textasciicircum{}0\_\allowbreak{}0 == f\_\allowbreak{}W (energy density) and T\textasciicircum{}1\_\allowbreak{}1 == f\_\allowbreak{}W (pressure -f\_\allowbreak{}W: w = -1 along the brane)} & True & 25.8 \\
26 & \texttt{sources smeared over the torus (W(a, b) the comoving energy, section 20.2): the brane tension W = a\textasciicircum{}3 f\_\allowbreak{}W/\allowbreak{}Lc\textasciicircum{}4 gives eps = f\_\allowbreak{}W/\allowbreak{}V\_\allowbreak{}perp, pa = -eps, pb = 0; brane radiation W = a\textasciicircum{}3 r0 a\textasciicircum{}-4/\allowbreak{}Lc\textasciicircum{}4 gives pa = eps/\allowbreak{}3, pb = 0; radiation of the bulk, isotropic in the seven space directions, has <cos\textasciicircum{}2 theta> = 3/\allowbreak{}7 and <sin\textasciicircum{}2 theta> = 4/\allowbreak{}7 for the weight cos\textasciicircum{}2 theta sin\textasciicircum{}3 theta of the momentum sphere, so pa = (1/\allowbreak{}3)(3/\allowbreak{}7) eps = eps/\allowbreak{}7 and pb = (1/\allowbreak{}4)(4/\allowbreak{}7) eps = eps/\allowbreak{}7} & True & 25.8 \\
27 & \texttt{the brane in four-dimensional physics: (i) the tension is vacuum energy, kappa f\_\allowbreak{}W/\allowbreak{}V\_\allowbreak{}perp == f\_\allowbreak{}W/\allowbreak{}M\_\allowbreak{}Pl\textasciicircum{}2 (V\_\allowbreak{}perp = kappa M\_\allowbreak{}Pl\textasciicircum{}2), and f\_\allowbreak{}W = M8\textasciicircum{}4 is between 3 10\textasciicircum{}62 and 5 10\textasciicircum{}62 times rho\_\allowbreak{}Lambda = 3 H0\textasciicircum{}2 M\_\allowbreak{}Pl\textasciicircum{}2 Omega\_\allowbreak{}Lambda (section 24.5); (ii) normalcy temperature: (T\textasciicircum{}7/\allowbreak{}M8\textasciicircum{}6)/\allowbreak{}(T\textasciicircum{}2/\allowbreak{}M\_\allowbreak{}Pl) == 1 at T* = (M8\textasciicircum{}6/\allowbreak{}M\_\allowbreak{}Pl)\textasciicircum{}(1/\allowbreak{}5) = 13.3 GeV, and V\_\allowbreak{}perp T\textasciicircum{}4 T\textasciicircum{}3/\allowbreak{}M\_\allowbreak{}Pl\textasciicircum{}2 == T\textasciicircum{}7/\allowbreak{}M8\textasciicircum{}6; (iii) thin brane: (1/\allowbreak{}y\textasciicircum{}3) d/\allowbreak{}dy (y\textasciicircum{}3 d/\allowbreak{}dy y\textasciicircum{}-2) == 0 and 2 Pi\textasciicircum{}2 y\textasciicircum{}3 d/\allowbreak{}dy y\textasciicircum{}-2 == -4 Pi\textasciicircum{}2 (the Green's function -1/\allowbreak{}(4 Pi\textasciicircum{}2 |y|\textasciicircum{}2) of the four transverse directions), and Sqrt[kappa f\_\allowbreak{}W] == 1/\allowbreak{}M8 for f\_\allowbreak{}W = M8\textasciicircum{}4} & True & 25.8 \\*
\end{longtable}}
The Python program’s boolean checks:

{\footnotesize
\begin{longtable}{>{\raggedright\arraybackslash}p{\dimexpr 0.0500\linewidth-2\tabcolsep\relax}>{\raggedright\arraybackslash}p{\dimexpr 0.8300\linewidth-2\tabcolsep\relax}>{\raggedright\arraybackslash}p{\dimexpr 0.1200\linewidth-2\tabcolsep\relax}}
\toprule
\textbf{\#} & \textbf{check} & \textbf{result} \\
\midrule\endfirsthead
\toprule
\textbf{\#} & \textbf{check} & \textbf{result} \\
\midrule\endhead
\endfoot
\bottomrule\endlastfoot
1 & \texttt{torus: det(-g\_\allowbreak{}kl) = b\textasciicircum{}8 for the cosmological metric; the proper volume of [0, Lc)\textasciicircum{}4 is b\textasciicircum{}4 Lc\textasciicircum{}4} & True \\
2 & \texttt{only b Lc is physical: y = lambda y' maps (Lc, b) to (Lc/\allowbreak{}lambda, lambda b); b\textasciicircum{}2 dy.dy, b Lc and b\textasciicircum{}4 Lc\textasciicircum{}4 unchanged} & True \\
3 & \texttt{symmetries fixing y = 0: 384 = 2\textasciicircum{}4 4! distinct signed permutation matrices, orthogonal integer matrices; shells 1, 2, 3, 4, 9 have orbits \{8\}, \{24\}, \{32\}, \{8, 16\}, \{8, 96\}} & True \\
4 & \texttt{Fourier modes phi\_\allowbreak{}n = exp(i q\_\allowbreak{}n.y)/\allowbreak{}sqrt(V\_\allowbreak{}perp), q\_\allowbreak{}n = 2 pi n/\allowbreak{}Lc: orthonormal for b\textasciicircum{}4 d\textasciicircum{}4y (three exact pairs)} & True \\
5 & \texttt{Parseval: f = prod cos\textasciicircum{}2(pi y\_\allowbreak{}k/\allowbreak{}Lc) has 81 coefficients sqrt(V\_\allowbreak{}perp) prod w(n\_\allowbreak{}k) (3 checked by integration), Sum |c\_\allowbreak{}n|\textasciicircum{}2 = Int b\textasciicircum{}4 f\textasciicircum{}2 = (3/\allowbreak{}8)\textasciicircum{}4 V\_\allowbreak{}perp} & True \\
6 & \texttt{Laplacian of the torus: -(1/\allowbreak{}b\textasciicircum{}2) Sum d\textasciicircum{}2/\allowbreak{}dy\_\allowbreak{}k\textasciicircum{}2 phi\_\allowbreak{}n = (|q\_\allowbreak{}n|\textasciicircum{}2/\allowbreak{}b\textasciicircum{}2) phi\_\allowbreak{}n} & True \\
7 & \texttt{scalar Kaluza-Klein masses: (Box + m\textasciicircum{}2) exp(-iEx0 + ik.x + iq.y) = (-E\textasciicircum{}2 + |k|\textasciicircum{}2 + |q|\textasciicircum{}2/\allowbreak{}b0\textasciicircum{}2 + m\textasciicircum{}2) exp(...) on diag(1,-1,-1,-1,-b0\textasciicircum{}2,...)} & True \\
8 & \texttt{graviton on the torus: D(D-3)/\allowbreak{}2 = 20 (D = 8) = 2 + 8 + 10 (zero mode: graviton, 4 graviphotons, 10 moduli) = 5 + 9 + 6 (massive level)} & True \\
9 & \texttt{Weyl's law: N(10) = 49689, N(100) = 493490641; |N(R)/\allowbreak{}(pi\textasciicircum{}2 R\textasciicircum{}4/\allowbreak{}2) - 1| < 1e-2, 1e-3, 3e-5, 5e-6, 1e-6} & True \\
10 & \texttt{numbers today: V\_\allowbreak{}perp = 5.929e12 GeV\textasciicircum{}-4, l\_\allowbreak{}perp = 1560.4 GeV\textasciicircum{}-1 = 3.079e-13 m, 2 pi/\allowbreak{}l\_\allowbreak{}perp = 4.027 MeV; N(5 GeV) = 1.17e13 (lattice to 1e-5); N(M8) = M\_\allowbreak{}Pl\textasciicircum{}2/\allowbreak{}(32 pi\textasciicircum{}2 M8\textasciicircum{}2)} & True \\
11 & \texttt{from eight to four dimensions: for diag(1, -a\textasciicircum{}2, -a\textasciicircum{}2, -a\textasciicircum{}2, -b0\textasciicircum{}2, ..., -b0\textasciicircum{}2) (static torus) R = R4 = -6(Ha' + 2 Ha\textasciicircum{}2), the scalar curvature of the FRW metric diag(1, -a\textasciicircum{}2, -a\textasciicircum{}2, -a\textasciicircum{}2); sqrt|g| = a\textasciicircum{}3 b0\textasciicircum{}4} & True \\
12 & \texttt{quaternion subalgebra of the brane directions: e1 e2 = -e3, e2 e3 = -e1, e3 e1 = -e2 (e\_\allowbreak{}i e\_\allowbreak{}j = -delta\_\allowbreak{}ij - eps\_\allowbreak{}ijk e\_\allowbreak{}k), span\{e0, ..., e3\} closed under the octonion product, Re(e\_\allowbreak{}A e\_\allowbreak{}B) = diag(1, -1, -1, -1)} & True \\
13 & \texttt{Newton in 8D: Omega\_\allowbreak{}6 = 16 pi\textasciicircum{}3/\allowbreak{}15, (1/\allowbreak{}r\textasciicircum{}6) d/\allowbreak{}dr (r\textasciicircum{}6 d/\allowbreak{}dr r\textasciicircum{}-5) = 0, Phi = -(5/\allowbreak{}6) kappa M/\allowbreak{}(5 Omega\_\allowbreak{}6 r\textasciicircum{}5) = -5 kappa M/\allowbreak{}(32 pi\textasciicircum{}3 r\textasciicircum{}5); D = 4: -G M/\allowbreak{}r} & True \\
14 & \texttt{4D Fourier transform: Int d\textasciicircum{}4y exp(-ip.y) (r\textasciicircum{}2 + |y|\textasciicircum{}2)\textasciicircum{}(-5/\allowbreak{}2) = (4 pi\textasciicircum{}2/\allowbreak{}(3 r)) exp(-p r) (mpmath, three (p, r), to 1e-20)} & True \\
15 & \texttt{Poisson duality: Sum\_\allowbreak{}n (r\textasciicircum{}2 + |n|\textasciicircum{}2)\textasciicircum{}(-5/\allowbreak{}2) = (4 pi\textasciicircum{}2/\allowbreak{}(3 r)) Sum\_\allowbreak{}m exp(-2 pi r |m|) to 1e-8 at r = 1/\allowbreak{}4, 1/\allowbreak{}2, 1, 2} & True \\
16 & \texttt{Z(5/\allowbreak{}2) = 7 zeta(5/\allowbreak{}2) zeta(3/\allowbreak{}2) = Sum r\_\allowbreak{}4(k) k\textasciicircum{}-5/\allowbreak{}2 (to 1e-8); near field ratio -> 1 (0.98 < . < 1 at r = 1/\allowbreak{}10, 1/\allowbreak{}20); far field ratio -> 1 (to 2e-2 at r = 2, 3)} & True \\
17 & \texttt{graviton + radion: (5 kappa M/\allowbreak{}(32 pi\textasciicircum{}3))(4 pi\textasciicircum{}2/\allowbreak{}(3 V\_\allowbreak{}perp r)) = (5/\allowbreak{}3) G4 M/\allowbreak{}r; radion (1/\allowbreak{}(sqrt3 M\_\allowbreak{}Pl))\textasciicircum{}2/\allowbreak{}(4 pi) = (2/\allowbreak{}3)/\allowbreak{}(8 pi M\_\allowbreak{}Pl\textasciicircum{}2)} & True \\
18 & \texttt{Kaluza-Klein momentum: bulk vertex Int b\textasciicircum{}4 phi\_\allowbreak{}n phi\_\allowbreak{}m phi\_\allowbreak{}l\textasciicircum{}* = 1/\allowbreak{}sqrt(V\_\allowbreak{}perp) for l = n + m, 0 otherwise; at y = 0 phi\_\allowbreak{}n phi\_\allowbreak{}m phi\_\allowbreak{}l\textasciicircum{}* = V\_\allowbreak{}perp\textasciicircum{}(-3/\allowbreak{}2)} & True \\
19 & \texttt{delta function of the torus: Fourier coefficients Int exp(-i q\_\allowbreak{}n y) delta(y) dy = 1; Int b\textasciicircum{}4 d\textasciicircum{}4y rho4 delta\textasciicircum{}4(y)/\allowbreak{}b\textasciicircum{}4 = rho4 (average rho4/\allowbreak{}V\_\allowbreak{}perp)} & True \\
20 & \texttt{induced metric h = J\textasciicircum{}T g J: FRW diag(1, -a\textasciicircum{}2, -a\textasciicircum{}2, -a\textasciicircum{}2) for the cosmological metric; diag(n\textasciicircum{}2, -a\textasciicircum{}2, -a\textasciicircum{}2, -a\textasciicircum{}2) at x4 = 0 with n = 1 + x0\textasciicircum{}2/\allowbreak{}11, a = 1 + x0/\allowbreak{}3 for the inhomogeneous metric} & True \\
21 & \texttt{brane volume element: sqrt|g| = sqrt|h| b\textasciicircum{}4 at the brane (cosmological metric; inhomogeneous metric at x0 = 1/\allowbreak{}2, exact), so Int\_\allowbreak{}W d\textasciicircum{}4x sqrt|h| L = Int d\textasciicircum{}8X sqrt|g| delta\textasciicircum{}4(y)/\allowbreak{}b\textasciicircum{}4 L; test Int b\textasciicircum{}4 (delta\textasciicircum{}4/\allowbreak{}b\textasciicircum{}4) F = F(0)} & True \\
22 & \texttt{Christoffel symbols \{k mu nu\} at the brane: all 0 for the cosmological metric (totally geodesic); for the inhomogeneous metric at x0 = 1/\allowbreak{}2 only \{4 00\} = 1458/\allowbreak{}58619 = n d4n/\allowbreak{}b\textasciicircum{}2} & True \\
23 & \texttt{R\_\allowbreak{}A = (gamma\_\allowbreak{}mu (x) 1\_\allowbreak{}4, gamma5 (x) rho\_\allowbreak{}k) and Gamma17[A] both satisfy X\_\allowbreak{}A X\_\allowbreak{}B + X\_\allowbreak{}B X\_\allowbreak{}A = 2 eta\_\allowbreak{}AB 1\_\allowbreak{}16} & True \\
24 & \texttt{every Gamma17[A] is a signed permutation matrix} & True \\
25 & \texttt{intertwiner: Schur average with the seed matrix unit (1, 2); 2 sqrt(2) U has entries 0, +-1, +-i (rounding error < 1e-12) with 8 non-zero entries in each row and column; Gamma17[A].U8 = U8.R\_\allowbreak{}A and U8\textasciicircum{}dag.U8 = 8 1\_\allowbreak{}16 exactly} & True \\
26 & \texttt{Schur's lemma: the linear map X -> (Gamma\_\allowbreak{}A X - X R\_\allowbreak{}A)\_\allowbreak{}A has a one-dimensional null space (one singular value < 1e-10, the next > 0.1)} & True \\
27 & \texttt{brane Lorentz group on the brane spinor index, transverse rotations on a: Gamma\_\allowbreak{}mu Gamma\_\allowbreak{}nu = U (gamma\_\allowbreak{}mu gamma\_\allowbreak{}nu (x) 1) U\textasciicircum{}dag, Gamma\_\allowbreak{}k Gamma\_\allowbreak{}l = U (1 (x) rho\_\allowbreak{}k rho\_\allowbreak{}l) U\textasciicircum{}dag, Gamma\_\allowbreak{}mu Gamma\_\allowbreak{}k = U (gamma\_\allowbreak{}mu gamma5 (x) rho\_\allowbreak{}k) U\textasciicircum{}dag} & True \\
28 & \texttt{Spin(4) on the transverse index: rho4567 = -gamma5 (4x4) commutes with rho\_\allowbreak{}k rho\_\allowbreak{}l, anticommutes with rho\_\allowbreak{}k; (rho\_\allowbreak{}k rho\_\allowbreak{}l + s rho\_\allowbreak{}m rho\_\allowbreak{}n) P\_\allowbreak{}s = 0 and (rho\_\allowbreak{}k rho\_\allowbreak{}l - s rho\_\allowbreak{}m rho\_\allowbreak{}n) P\_\allowbreak{}s has rank 2 for even (k,l,m,n): 4 = (2,1) + (1,2)} & True \\
29 & \texttt{chirality: chi = i Gamma\_\allowbreak{}0...Gamma\_\allowbreak{}7 = U (gamma5 (x) rho4567) U\textasciicircum{}dag; P\_\allowbreak{}L17 = U (P\_\allowbreak{}L (x) P\_\allowbreak{}+ + P\_\allowbreak{}R (x) P\_\allowbreak{}-) U\textasciicircum{}dag, each part of rank 4} & True \\
30 & \texttt{example psi = U.(u (x) e\_\allowbreak{}1), u = (1, 2, 0, -1): psibar psi = ubar u = 4 and psibar Gamma\textasciicircum{}mu psi = ubar gamma\textasciicircum{}mu u} & True \\
31 & \texttt{four brane Dirac spinors: psibar psi' = Sum\_\allowbreak{}a psibar\_\allowbreak{}a psi'\_\allowbreak{}a and psibar Gamma\textasciicircum{}mu psi' = Sum\_\allowbreak{}a psibar\_\allowbreak{}a gamma\textasciicircum{}mu psi'\_\allowbreak{}a (random integer spinors)} & True \\
32 & \texttt{Kaluza-Klein operator X = gamma5 (x) (rho.q)/\allowbreak{}b0: X\textasciicircum{}2 = -|q|\textasciicircum{}2/\allowbreak{}b0\textasciicircum{}2 1\_\allowbreak{}16 and X anticommutes with gamma\_\allowbreak{}mu (x) 1\_\allowbreak{}4} & True \\
33 & \texttt{brane tension: T\_\allowbreak{}mu nu = (2/\allowbreak{}sqrt|h|) dS/\allowbreak{}dh\textasciicircum{}mu nu = f\_\allowbreak{}W h\_\allowbreak{}mu nu for S = -Int sqrt|h| f\_\allowbreak{}W; T\textasciicircum{}0\_\allowbreak{}0 = f\_\allowbreak{}W, T\textasciicircum{}1\_\allowbreak{}1 = f\_\allowbreak{}W (p = -f\_\allowbreak{}W, w = -1)} & True \\
34 & \texttt{smeared sources: tension eps = f\_\allowbreak{}W/\allowbreak{}V\_\allowbreak{}perp, pa = -eps, pb = 0; brane radiation pa = eps/\allowbreak{}3, pb = 0; bulk radiation <cos\textasciicircum{}2> = 3/\allowbreak{}7, <sin\textasciicircum{}2> = 4/\allowbreak{}7: pa = pb = eps/\allowbreak{}7} & True \\
35 & \texttt{normalcy temperature: (T\textasciicircum{}7/\allowbreak{}M8\textasciicircum{}6)/\allowbreak{}(T\textasciicircum{}2/\allowbreak{}M\_\allowbreak{}Pl) = 1 at T* = (M8\textasciicircum{}6/\allowbreak{}M\_\allowbreak{}Pl)\textasciicircum{}(1/\allowbreak{}5) = 13.3 GeV; V\_\allowbreak{}perp T\textasciicircum{}4 T\textasciicircum{}3/\allowbreak{}M\_\allowbreak{}Pl\textasciicircum{}2 = T\textasciicircum{}7/\allowbreak{}M8\textasciicircum{}6} & True \\
36 & \texttt{the tension is vacuum energy: kappa f\_\allowbreak{}W/\allowbreak{}V\_\allowbreak{}perp = f\_\allowbreak{}W/\allowbreak{}M\_\allowbreak{}Pl\textasciicircum{}2 (V\_\allowbreak{}perp = kappa M\_\allowbreak{}Pl\textasciicircum{}2); f\_\allowbreak{}W = M8\textasciicircum{}4 is 3e62 to 5e62 rho\_\allowbreak{}Lambda} & True \\
37 & \texttt{Green's function of the four transverse directions: (1/\allowbreak{}y\textasciicircum{}3) d/\allowbreak{}dy (y\textasciicircum{}3 d/\allowbreak{}dy y\textasciicircum{}-2) = 0, flux 2 pi\textasciicircum{}2 y\textasciicircum{}3 d/\allowbreak{}dy y\textasciicircum{}-2 = -4 pi\textasciicircum{}2; sqrt(kappa f\_\allowbreak{}W) = 1/\allowbreak{}M8 for f\_\allowbreak{}W = M8\textasciicircum{}4} & True \\*
\end{longtable}}
The numbers printed by both programs (largest relative difference per group; integers must agree exactly):

{\footnotesize
\begin{longtable}{>{\raggedright\arraybackslash}p{\dimexpr 0.5200\linewidth-2\tabcolsep\relax}>{\raggedright\arraybackslash}p{\dimexpr 0.1000\linewidth-2\tabcolsep\relax}>{\raggedright\arraybackslash}p{\dimexpr 0.2200\linewidth-2\tabcolsep\relax}>{\raggedright\arraybackslash}p{\dimexpr 0.1600\linewidth-2\tabcolsep\relax}}
\toprule
\textbf{quantity} & \textbf{values} & \textbf{max. relative difference} & \textbf{tolerance} \\
\midrule\endfirsthead
\toprule
\textbf{quantity} & \textbf{values} & \textbf{max. relative difference} & \textbf{tolerance} \\
\midrule\endhead
\endfoot
\bottomrule\endlastfoot
orbit sizes of the shells $k=1,\dots,12$ & 16 & 0 & exact \\
$r_4(k)$, $k=0,\dots,16$ & 17 & 0 & exact \\
Weyl: $N(R)$, $R=10,30,100,300,1000$ & 5 & 0 & exact \\
Weyl: $N(R)/(\pi^2R^4/2)-1$ & 5 & $3.8\times10^{-10}$ & $1\times10^{-9}$ \\
$V_\perp$, $\ell_\perp$ (GeV$^{-1}$, m, fm), $2\pi/\ell_\perp$ & 5 & $2.2\times10^{-16}$ & $1\times10^{-12}$ \\
$\ell_\perp/r_p$, $a_0/\ell_\perp$, $R=\ell_\perp/2\pi$, $M_D$ & 5 & $1.7\times10^{-16}$ & $1\times10^{-12}$ \\
$N(\mu)$ at 5 GeV, $m_h$, $M_8$; lattice count below 5 GeV & 4 & $1.8\times10^{-16}$ & $1\times10^{-12}$ \\
the first levels ($k$, $r_4$, $4r_4$, mass) & 32 & $2.5\times10^{-16}$ & $1\times10^{-12}$ \\
image sums $S(r)$ and Kaluza–Klein sums, $r=\frac14,\frac12,1,2$ & 8 & $1.2\times10^{-15}$ & $1\times10^{-9}$ \\
$Z(5/2)$: closed form, shell sum & 2 & $2.9\times10^{-16}$ & $1\times10^{-9}$ \\
near-field and far-field ratios & 5 & $7.8\times10^{-12}$ & $1\times10^{-6}$ \\
Christoffel symbol $\{4\,00\}$ at the brane & 1 & 0 & $1\times10^{-15}$ \\
$\Gamma_A$ as signed permutations & 128 & 0 & exact \\
$\mathrm{Re}(2\sqrt2\,U)$ & 256 & 0 & exact \\
$\mathrm{Im}(2\sqrt2\,U)$ & 256 & 0 & exact \\
example spinor: $\bar\psi\psi$, $\bar\psi\Gamma^\mu\psi$ & 5 & 0 & $1\times10^{-15}$ \\
$T_*$, $f_W/\rho_\Lambda$, $1/(2\pi M_8)$ & 3 & $5.8\times10^{-16}$ & $1\times10^{-12}$ \\*
\end{longtable}}
(Checks 1–14 concern the torus alone, sections 25.1–25.4; checks 15–16 the torus and the brane together, section 25.5; checks 17–27 the world brane, sections 25.6–25.8.) The results quoted from the other sections of 7+1.nb rest on these checks of the notebook, all evaluated in the same run:

{\footnotesize
\begin{longtable}{>{\raggedright\arraybackslash}p{\dimexpr 0.0500\linewidth-2\tabcolsep\relax}>{\raggedright\arraybackslash}p{\dimexpr 0.7700\linewidth-2\tabcolsep\relax}>{\raggedright\arraybackslash}p{\dimexpr 0.0800\linewidth-2\tabcolsep\relax}>{\raggedright\arraybackslash}p{\dimexpr 0.1000\linewidth-2\tabcolsep\relax}}
\toprule
\textbf{\#} & \textbf{check} & \textbf{result} & \textbf{section} \\
\midrule\endfirsthead
\toprule
\textbf{\#} & \textbf{check} & \textbf{result} & \textbf{section} \\
\midrule\endhead
\endfoot
\bottomrule\endlastfoot
1 & \texttt{scalar curvature R = -2 (3Ha' + 6Ha\textasciicircum{}2 + 4Hb' + 10Hb\textasciicircum{}2 + 12 Ha Hb)} & True & 11.2 \\
2 & \texttt{Einstein Kasner: (p, q) = ((1+2Sqrt[2])/\allowbreak{}7, (2-3Sqrt[2])/\allowbreak{}14) and ((1-2Sqrt[2])/\allowbreak{}7, (2+3Sqrt[2])/\allowbreak{}14) solve E\textasciicircum{}(1) == 0} & True & 14.1 \\
3 & \texttt{Gauss-Bonnet Kasner: 3 real roots of the cubic; (p, (3-3p)/\allowbreak{}4) solve E\textasciicircum{}(2) == 0} & True & 14.1 \\
4 & \texttt{static auxiliary space (Hb = 0): E\textasciicircum{}(1) = (3Ha\textasciicircum{}2, 2Ha' + 3Ha\textasciicircum{}2, 3Ha' + 6Ha\textasciicircum{}2), E\textasciicircum{}(2) = (0, 0, -12 Ha\textasciicircum{}2 (Ha' + Ha\textasciicircum{}2)), E\textasciicircum{}(3) = 0} & True & 14.4 \\
5 & \texttt{ADD relation: kappa8 rho4/\allowbreak{}V\_\allowbreak{}perp == 8 Pi G4 rho4 with G4 = G8/\allowbreak{}V\_\allowbreak{}perp (kappa8 = 8 Pi G8)} & True & 17.5 \\
6 & \texttt{KK operator: X.X == -|q|\textasciicircum{}2/\allowbreak{}b0\textasciicircum{}2 ID16 and X anticommutes with gamma\_\allowbreak{}mu (x) 1\_\allowbreak{}4} & True & 18.1 \\
7 & \texttt{KK masses: Det[gamma.p (x) 1 + X - m] vanishes on p0\textasciicircum{}2 - |p|\textasciicircum{}2 == m\textasciicircum{}2 + |q|\textasciicircum{}2/\allowbreak{}b0\textasciicircum{}2 (40-digit arithmetic) and not off shell} & True & 18.1 \\
8 & \texttt{stabilizing potential: W = a\textasciicircum{}3 b\textasciicircum{}4 U(b) gives eps = U, pa = -U, pb = -U - (b/\allowbreak{}4) U'(b)} & True & 20.2 \\
9 & \texttt{brane tension smeared over the torus (V\_\allowbreak{}perp = b\textasciicircum{}4 Lc\textasciicircum{}4): eps = f\_\allowbreak{}W/\allowbreak{}V\_\allowbreak{}perp, pa = -f\_\allowbreak{}W/\allowbreak{}V\_\allowbreak{}perp, pb = 0} & True & 20.2 \\
10 & \texttt{brane fluid rho\_\allowbreak{}4 = N a\textasciicircum{}(-3(1+w)), smeared: eps = rho\_\allowbreak{}4/\allowbreak{}V\_\allowbreak{}perp, pa = w eps, pb = 0} & True & 20.2 \\
11 & \texttt{radion equation of Einstein gravity: Hb' + 3 Ha Hb + 4 Hb\textasciicircum{}2 == kappa (eps/\allowbreak{}6 - pa/\allowbreak{}2 + pb/\allowbreak{}3) + Lambda/\allowbreak{}3 (constraint used)} & True & 20.3 \\
12 & \texttt{radion mass: U = U0 + (K\_\allowbreak{}U/\allowbreak{}2) beta\textasciicircum{}2, b U' = K\_\allowbreak{}U beta gives the restoring term -(kappa K\_\allowbreak{}U/\allowbreak{}12) beta, m\_\allowbreak{}r\textasciicircum{}2 = kappa K\_\allowbreak{}U/\allowbreak{}12} & True & 20.3 \\
13 & \texttt{static auxiliary space: (kappa/\allowbreak{}4) b U' == Lambda + kappa U + 2 kappa f\_\allowbreak{}W/\allowbreak{}V\_\allowbreak{}perp + (kappa/\allowbreak{}2)(eps - 3pa) + kappa pb - 12 alpha2 Ha\textasciicircum{}2 (Ha' + Ha\textasciicircum{}2)} & True & 20.3 \\
14 & \texttt{brane Friedmann equation: 3 Ha\textasciicircum{}2 == (kappa/\allowbreak{}V\_\allowbreak{}perp) rho\_\allowbreak{}4 + Lambda4 with Lambda4 = Lambda + kappa U0 + kappa f\_\allowbreak{}W/\allowbreak{}V\_\allowbreak{}perp, i.e. 8 Pi G4 = kappa/\allowbreak{}V\_\allowbreak{}perp} & True & 20.3 \\
15 & \texttt{Einstein-frame form of the constraint: 3 (Ha + 2 Hb)\textasciicircum{}2 == 6 Hb\textasciicircum{}2 + (E\textasciicircum{}(1)0\_\allowbreak{}0 of section 12.6), i.e. 3 M\_\allowbreak{}Pl\textasciicircum{}2 H\_\allowbreak{}E\textasciicircum{}2 == phi'\textasciicircum{}2/\allowbreak{}2 + rho with phi = Sqrt[12] M\_\allowbreak{}Pl beta} & True & 20.3 \\
16 & \texttt{Kasner constants: P\_\allowbreak{}E(u1) == 0, p1 == (1 + 2 Sqrt[2])/\allowbreak{}7, q1 == (2 - 3 Sqrt[2])/\allowbreak{}14; u2 in (-1/\allowbreak{}3, 0) with P\_\allowbreak{}G(u2) == 0, 3 p2 + 4 q2 == 3} & True & 21.1 \\
17 & \texttt{cubic regime: Ha = 1/\allowbreak{}tau + A0 + A1 tau, Hb = h0 + h1 tau + h2 tau\textasciicircum{}2 solves the 00, 11, 44 equations through order tau\textasciicircum{}-1 (11: checked, not used)} & True & 21.2 \\
18 & \texttt{matter-dominated cubic regime: p\_\allowbreak{}eff -> 3/\allowbreak{}4, Hb -> -Sqrt[alpha2/\allowbreak{}(54 alpha3)], kappa eps tau\textasciicircum{}3 -> 9 alpha2 |h*|, alpha2 E2 : alpha3 E3 : kappa eps -> 9/\allowbreak{}4 : -5/\allowbreak{}4 : 1 (Ha = 10\textasciicircum{}15)} & True & 21.3 \\
19 & \texttt{without the stabilizing potential the gas-dominated Einstein regime decompactifies: tau Ha -> 0, tau Hb -> 2/\allowbreak{}5, pb/\allowbreak{}eps -> 1/\allowbreak{}4 (Ha = 10\textasciicircum{}-14)} & True & 21.3 \\
20 & \texttt{Einstein + fluid power laws a \textasciitilde{} tau\textasciicircum{}P, b \textasciitilde{} tau\textasciicircum{}Q: (wa, wb) = (0, 1/\allowbreak{}4) -> (0, 2/\allowbreak{}5); (1/\allowbreak{}3, 0) -> (1/\allowbreak{}2, 0); (1/\allowbreak{}7, 1/\allowbreak{}7) -> (1/\allowbreak{}4, 1/\allowbreak{}4) (tau\textasciicircum{}2 E\textasciicircum{}(1) from EKas of section 14.1)} & True & 21.3 \\
21 & \texttt{capture: the radion oscillates like a four-dimensional massive field in a matter-dominated universe: with tau\_\allowbreak{}E from <H\_\allowbreak{}E> = 2/\allowbreak{}(3 (tau - tau\_\allowbreak{}E)) over the last period, beta = A sin(m\_\allowbreak{}r tau + vartheta)/\allowbreak{}(m\_\allowbreak{}r (tau - tau\_\allowbreak{}E)) to 2 10\textasciicircum{}-2, <H\_\allowbreak{}E>(tau - tau\_\allowbreak{}E) = 2/\allowbreak{}3 to 10\textasciicircum{}-2 at tau = 10\textasciicircum{}5 and 2 10\textasciicircum{}5, the gas share constant to 2 10\textasciicircum{}-2} & True & 22.1 \\
22 & \texttt{the number of lattice vectors with |n|\textasciicircum{}2 = k is r\_\allowbreak{}4(k) = 8 sigma(k) - 32 sigma(k/\allowbreak{}4) (direct count for k <= 60)} & True & 22.2 \\
23 & \texttt{the radion decays into the bulk: the lattice sum over the Kaluza-Klein modes of dirac16complex (m << m\_\allowbreak{}r/\allowbreak{}2) equals its continuum limit to 10\textasciicircum{}-6, the number of momenta to 10\textasciicircum{}-5, and Gamma\_\allowbreak{}bulk/\allowbreak{}Gamma\_\allowbreak{}brane > 10\textasciicircum{}12} & True & 22.2 \\
24 & \texttt{radion equilibrium in de Sitter: 3 m\_\allowbreak{}r\textasciicircum{}2 beta\_\allowbreak{}eq == 3 H\textasciicircum{}2 - 12 alpha2 H\textasciicircum{}4 + O(beta\_\allowbreak{}eq\textasciicircum{}2)} & True & 23.2 \\
25 & \texttt{eight-dimensional de Sitter (no stabilizer): all three equations read 21 H\textasciicircum{}2 - 420 alpha2 H\textasciicircum{}4 + 2520 alpha3 H\textasciicircum{}6 == Lambda} & True & 23.2 \\
26 & \texttt{Kaluza-Klein tower: Sum\_\allowbreak{}n (1 - mu\_\allowbreak{}n\textasciicircum{}2/\allowbreak{}m\_\allowbreak{}h\textasciicircum{}2)\textasciicircum{}2 -> V\_\allowbreak{}perp m\_\allowbreak{}h\textasciicircum{}4/\allowbreak{}(192 Pi\textasciicircum{}2) and Sum\_\allowbreak{}n mu\_\allowbreak{}n (1 - mu\_\allowbreak{}n\textasciicircum{}2/\allowbreak{}m\_\allowbreak{}h\textasciicircum{}2)\textasciicircum{}2 -> V\_\allowbreak{}perp m\_\allowbreak{}h\textasciicircum{}5/\allowbreak{}(315 Pi\textasciicircum{}2) (d\textasciicircum{}4n = 2 Pi\textasciicircum{}2 |n|\textasciicircum{}3 d|n|: exact integrals); the lattice sums with r\_\allowbreak{}4(k) for R = 100 equal Pi\textasciicircum{}2 R\textasciicircum{}4/\allowbreak{}12 and 16 Pi\textasciicircum{}2 R\textasciicircum{}4/\allowbreak{}315 to 10\textasciicircum{}-8} & True & 24.4 \\
27 & \texttt{Epstein sum over Z\textasciicircum{}4 without 0: Sum |n|\textasciicircum{}-8 == 8 (1 - 4\textasciicircum{}-3) zeta(4) zeta(3), compared with the shell sum Sum r\_\allowbreak{}4(k)/\allowbreak{}k\textasciicircum{}4 to k = 20000 (tail < 3 10\textasciicircum{}-8)} & True & 24.5 \\
28 & \texttt{Casimir energy of dirac16complex on the auxiliary torus: massless 16 Gamma(4) Z/\allowbreak{}(Pi\textasciicircum{}4 l\textasciicircum{}4) > 10\textasciicircum{}34 rho\_\allowbreak{}Lambda, the massive sum -> the massless value for m l -> 0 (to 10\textasciicircum{}-4), and equals rho\_\allowbreak{}Lambda at m l = 92.2} & True & 24.5 \\*
\end{longtable}}
Section 10 rests on these checks of section 26 of 7+1.nb (which contains 33 checks, 33 True; the independent program tools/ref\_radion.py has 31 boolean checks, 31 True, and the 15 groups of numbers printed by both agree to a largest relative difference of $8.2\times10^{-11}$):

{\footnotesize
\begin{longtable}{>{\raggedright\arraybackslash}p{\dimexpr 0.0500\linewidth-2\tabcolsep\relax}>{\raggedright\arraybackslash}p{\dimexpr 0.7700\linewidth-2\tabcolsep\relax}>{\raggedright\arraybackslash}p{\dimexpr 0.0800\linewidth-2\tabcolsep\relax}>{\raggedright\arraybackslash}p{\dimexpr 0.1000\linewidth-2\tabcolsep\relax}}
\toprule
\textbf{\#} & \textbf{check} & \textbf{result} & \textbf{section} \\
\midrule\endfirsthead
\toprule
\textbf{\#} & \textbf{check} & \textbf{result} & \textbf{section} \\
\midrule\endhead
\endfoot
\bottomrule\endlastfoot
1 & \texttt{the 4-sphere of radius r0: Ricci\_\allowbreak{}ij = 3 gamma\_\allowbreak{}ij (gamma the unit-sphere metric), Ricci\_\allowbreak{}ij = 0 for i != j, so R(S\textasciicircum{}4) = 12/\allowbreak{}r0\textasciicircum{}2 (in the signature of 7+1.nb the space-like block gives R = -12/\allowbreak{}r0\textasciicircum{}2, also with curvature[] of section 11); the volume of the unit S\textasciicircum{}4 is 8 Pi\textasciicircum{}2/\allowbreak{}3} & True & 26.1 \\
2 & \texttt{radion on T\textasciicircum{}4: for ds\textasciicircum{}2 = e\textasciicircum{}(-4 beta) (N\textasciicircum{}2 dx0\textasciicircum{}2 - a\textasciicircum{}2 dx.dx) - e\textasciicircum{}(2 beta) b0\textasciicircum{}2 dy.dy the reduced action -(1/\allowbreak{}(2 kappa)) Int(R + 2 Lambda) - Int U equals -(M\_\allowbreak{}Pl\textasciicircum{}2/\allowbreak{}2) R\_\allowbreak{}E + 6 M\_\allowbreak{}Pl\textasciicircum{}2 (d beta)\textasciicircum{}2 - e\textasciicircum{}(-4 beta)(M\_\allowbreak{}Pl\textasciicircum{}2 Lambda + V0 U) up to a total derivative (Euler-Lagrange expressions for N, a, beta vanish), M\_\allowbreak{}Pl\textasciicircum{}2 = V0/\allowbreak{}kappa, V0 = b0\textasciicircum{}4 Lc\textasciicircum{}4} & True & 26.1 \\
3 & \texttt{radion on S\textasciicircum{}4: the same reduction for ds\textasciicircum{}2 = e\textasciicircum{}(-4 beta) g\_\allowbreak{}E - e\textasciicircum{}(2 beta) r0\textasciicircum{}2 gamma (unit-sphere gamma) holds with the additional potential -(M\_\allowbreak{}Pl\textasciicircum{}2/\allowbreak{}2) e\textasciicircum{}(-4 beta) 12/\allowbreak{}(r0\textasciicircum{}2 e\textasciicircum{}(2 beta)) = -6 M\_\allowbreak{}Pl\textasciicircum{}2 e\textasciicircum{}(-6 beta)/\allowbreak{}r0\textasciicircum{}2 from the curvature of the sphere, V0 = (8 Pi\textasciicircum{}2/\allowbreak{}3) r0\textasciicircum{}4; R is independent of the angles} & True & 26.1 \\
4 & \texttt{radion on T\textasciicircum{}4 with U = U0 + U1 beta + (K\_\allowbreak{}U/\allowbreak{}2) beta\textasciicircum{}2: a flat vacuum at b = b0 needs Lambda + kappa U0 = 0 and U1 = 0, and m\_\allowbreak{}r\textasciicircum{}2 = V''(beta)/\allowbreak{}(12 M\_\allowbreak{}Pl\textasciicircum{}2) = kappa K\_\allowbreak{}U/\allowbreak{}12 (section 20.3)} & True & 26.2 \\
5 & \texttt{radion on S\textasciicircum{}4 with U: Lambda + kappa U0 == 6/\allowbreak{}r0\textasciicircum{}2, kappa U1 == -12/\allowbreak{}r0\textasciicircum{}2 and m\_\allowbreak{}r\textasciicircum{}2 == kappa K\_\allowbreak{}U/\allowbreak{}12 - 2/\allowbreak{}r0\textasciicircum{}2 from V(beta); the eight-dimensional radion equation gives the same static solution and the same mass} & True & 26.2 \\
6 & \texttt{S\textasciicircum{}4 with Lambda only: dS4 x S\textasciicircum{}4 with r0\textasciicircum{}2 = 9/\allowbreak{}Lambda, H\textasciicircum{}2 = 1/\allowbreak{}r0\textasciicircum{}2 = Lambda/\allowbreak{}9 solves the 00, 11 and 44 equations; it is a MAXIMUM of V: m\textasciicircum{}2 = -6 H\textasciicircum{}2 from V(beta), and m\textasciicircum{}2 = -6/\allowbreak{}r0\textasciicircum{}2 = -6 H\textasciicircum{}2 from the eight-dimensional radion equation} & True & 26.2 \\
7 & \texttt{zonal harmonics on S\textasciicircum{}4 (l = 0..6, exact): (1/\allowbreak{}sin\textasciicircum{}3 x4) d/\allowbreak{}dx4 (sin\textasciicircum{}3 x4 d/\allowbreak{}dx4 C\_\allowbreak{}l\textasciicircum{}(3/\allowbreak{}2)(cos x4)) == -l (l + 3) C\_\allowbreak{}l\textasciicircum{}(3/\allowbreak{}2), C\_\allowbreak{}l\textasciicircum{}(3/\allowbreak{}2)(1) == (l + 1)(l + 2)/\allowbreak{}2, and Int over S\textasciicircum{}4 of (C\_\allowbreak{}l/\allowbreak{}C\_\allowbreak{}l(1))\textasciicircum{}2 == V/\allowbreak{}d\_\allowbreak{}l (unit sphere), so the normalized zonal harmonic is Sqrt[d\_\allowbreak{}l/\allowbreak{}V] at the north pole; d\_\allowbreak{}l == Binomial[l + 4, 4] - Binomial[l + 2, 4] for l = 0..20} & True & 26.4 \\
8 & \texttt{Weyl's law on S\textasciicircum{}4: the number of levels with l(l + 3) <= (mu r)\textasciicircum{}2 approaches (mu r)\textasciicircum{}4/\allowbreak{}12 = V mu\textasciicircum{}4/\allowbreak{}(32 Pi\textasciicircum{}2) (V = 8 Pi\textasciicircum{}2 r\textasciicircum{}4/\allowbreak{}3): relative deviation -1.0 10\textasciicircum{}-2, -1.0 10\textasciicircum{}-4, -1.0 10\textasciicircum{}-6 at mu r = 10, 100, 1000 (like -1/\allowbreak{}(mu r)\textasciicircum{}2)} & True & 26.4 \\
9 & \texttt{the brane at the north pole of S\textasciicircum{}4: the embedding (r0 sin x4 sin x5 sin x6 cos x7, ..., r0 cos x4) of S\textasciicircum{}4 in R\textasciicircum{}5 induces r0\textasciicircum{}2 gamma; the north pole x4 = 0 is (0, 0, 0, 0, r0); of the 10 rotations L\_\allowbreak{}ab of R\textasciicircum{}5 (the Killing vectors of S\textasciicircum{}4, SO(5)) the 6 with a, b <= 4 vanish there (SO(4), kept by the brane) and the 4 L\_\allowbreak{}a5 move it} & True & 26.4 \\
10 & \texttt{Killing spinors of S\textasciicircum{}4 (radius r0, frame e\textasciicircum{}a = (r0, r0 sin x4, ...) dx\textasciicircum{}(a+3), gt\_\allowbreak{}a = I rho\_\allowbreak{}(a+3)): eps\_\allowbreak{}s = exp(s (i/\allowbreak{}2) x4 gt\_\allowbreak{}1) exp((1/\allowbreak{}2) x5 gt\_\allowbreak{}1 gt\_\allowbreak{}2) exp((1/\allowbreak{}2) x6 gt\_\allowbreak{}2 gt\_\allowbreak{}3) exp((1/\allowbreak{}2) x7 gt\_\allowbreak{}3 gt\_\allowbreak{}4) eps0 satisfies nabla\_\allowbreak{}mu eps\_\allowbreak{}s == s (i/\allowbreak{}(2 r0)) e\_\allowbreak{}mu\textasciicircum{}a gt\_\allowbreak{}a eps\_\allowbreak{}s (s = +-1, all eps0), so the Dirac operator gives (2 s i/\allowbreak{}r0) eps\_\allowbreak{}s: |lambda| = 2/\allowbreak{}r0, Friedrich's bound n R/\allowbreak{}(4(n - 1)) = 4/\allowbreak{}r0\textasciicircum{}2} & True & 26.4 \\
11 & \texttt{Dirac spectrum on S\textasciicircum{}4: Weyl's law for spinors (4 components) N(|lambda| <= L/\allowbreak{}r) \textasciitilde{} 4 V (L/\allowbreak{}r)\textasciicircum{}4/\allowbreak{}(32 Pi\textasciicircum{}2) = L\textasciicircum{}4/\allowbreak{}3 agrees with the degeneracies 2 x 4 Binomial[l + 3, 3] of the eigenvalues +-i(l + 2)/\allowbreak{}r (to 2.0 10\textasciicircum{}-3 at L = 1000)} & True & 26.4 \\
12 & \texttt{the sphere with the volume of the torus today (V = M\_\allowbreak{}Pl\textasciicircum{}2/\allowbreak{}M8\textasciicircum{}6): r = (3 V/\allowbreak{}(8 Pi\textasciicircum{}2))\textasciicircum{}(1/\allowbreak{}4) = 688.9 GeV\textasciicircum{}-1 = 1.359 10\textasciicircum{}-13 m; lowest scalar levels 4.03 MeV (torus, 8 states) and 2.90 MeV (sphere, l = 1, 5 states); lowest dirac16complex level on the sphere 2/\allowbreak{}r = 2.90 MeV (8 Dirac fermions; the torus has 4 massless ones for m = 0); levels below 5 GeV: torus 1.17335 10\textasciicircum{}13, sphere 1.17375 10\textasciicircum{}13, Weyl 1.17335 10\textasciicircum{}13} & True & 26.4 \\*
\end{longtable}}
\Needspace{12\baselineskip}\section*{13. How to reproduce}\addcontentsline{toc}{section}{13. How to reproduce}
\begin{enumerate}
\item Write the unevaluated notebook with \texttt{python3 tools/\allowbreak{}make\_\allowbreak{}7p1\_\allowbreak{}nb.py}; sections 25 and 26 are generated by \texttt{tools/\allowbreak{}nb\_\allowbreak{}part\_\allowbreak{}torus\_\allowbreak{}brane.py} and \texttt{tools/\allowbreak{}nb\_\allowbreak{}part\_\allowbreak{}radion.py}.
\item Evaluate 7+1.nb in a fresh kernel of Wolfram 15.0.1 (Evaluation \texorpdfstring{\ensuremath{\triangleright}}{▸} Quit Kernel, then Evaluation \texorpdfstring{\ensuremath{\triangleright}}{▸} Evaluate Notebook; or \texttt{Get["D:/\allowbreak{}Developer/\allowbreak{}github/\allowbreak{}7+1/\allowbreak{}tools/\allowbreak{}run\_\allowbreak{}eval.wl"]}) and save it.
\item Run \texttt{python3 tools/\allowbreak{}ref\_\allowbreak{}torus\_\allowbreak{}brane.py} and \texttt{python3 tools/\allowbreak{}ref\_\allowbreak{}radion.py}: they repeat the examples independently (exit code 0 if all pass).
\item Run \texttt{python3 tools/\allowbreak{}make\_\allowbreak{}docs.py}: it regenerates this document, its symbol dictionary and its provenance record (and all other documents).
\end{enumerate}
\Needspace{12\baselineskip}\section*{14. References}\addcontentsline{toc}{section}{14. References}
\begin{itemize}
\item Th. Kaluza, Sitzungsber. Preuss. Akad. Wiss. (1921) 966; O. Klein, Z. Phys. 37 (1926) 895 — a compact extra dimension.
\item N. Arkani-Hamed, S. Dimopoulos, G. Dvali, Phys. Lett. B 429 (1998) 263 and Phys. Rev. D 59 (1999) 086004 — large extra dimensions, $M_{\rm Pl}^2=V_\perp M_D^{n+2}$.
\item A. Kehagias, K. Sfetsos, Phys. Lett. B 472 (2000) 39 — the Yukawa corrections to Newton’s law from a torus.
\item C. G. J. Jacobi, Fundamenta nova theoriae functionum ellipticarum (1829) — the four-square theorem; H. Weyl, Math. Ann. 71 (1912) 441 — the asymptotic distribution of eigenvalues.
\item H. B. G. Casimir, Proc. K. Ned. Akad. Wet. 51 (1948) 793 — the Casimir energy; P. Epstein, Math. Ann. 56 (1903) 615 — lattice zeta functions.
\item S. Navas et al. (Particle Data Group), Phys. Rev. D 110 (2024) 030001 — searches for large extra dimensions.
\item R. Camporesi, A. Higuchi, J. Geom. Phys. 20 (1996) 1 — the Dirac spectrum on spheres; Th. Friedrich, Math. Nachr. 97 (1980) 117 — its lower bound.
\item E. Tiesinga, P. J. Mohr, D. B. Newell, B. N. Taylor, Rev. Mod. Phys. 93 (2021) 025010 (CODATA 2018) — $\hbar c$, $a_0$, $r_p$.
\end{itemize}
\end{document}
